\section{Introduction}
The analysis presented in this chapter is based on Ref.~\cite{Bhattacharyya:2025heterotic}.
{\color{black}
In this chapter we take a complementary approach to the classical two-body scattering problem, namely the scattering-amplitude approach. In the previous chapter, we saw that infrared divergences can arise in the computation of the eikonal phase within a bottom-up effective field theory, formulated in the worldline quantum field theory framework. A common expectation is that such infrared divergences cancel in physical observables once the dissipative contributions are included and the appropriate causal \(i0\) prescription is imposed.

The central question we address here is how robust this infrared finiteness really is. In particular, we would like to understand whether the cancellation of infrared divergences is a consequence of the specific formalism used, or whether it reflects a more general property of the underlying theory. The primary goal of this chapter is therefore to revisit the problem using amplitude methods, rather than the WQFT formalism, and to examine systematically how the infrared subtraction works in this complementary language. 

If, after a systematic subtraction of the known infrared contributions, residual infrared divergences still remain, there are two possible interpretations. One possibility is that the formalism or subtraction procedure is incomplete, and that additional conservative or radiative contributions have not been properly included. The other possibility is more fundamental: the persistence of such divergences may point to a genuine issue in the theory itself, or at least to a limitation in the way the classical observable is being defined. Thus, by comparing the amplitude-based analysis with the WQFT results, we aim to test the robustness of infrared finiteness and clarify whether the cancellation of infrared divergences is formalism-dependent or physically universal.  To sharpen this test, we consider a top-down effective field theory rather than a purely bottom-up model. In particular, we study Einstein--Maxwell--Dilaton theory, which arises as a low-energy effective theory of heterotic string theory. This choice provides a setting in which the interactions are not introduced arbitrarily, but are instead constrained by an ultraviolet completion. It therefore offers a useful arena to investigate whether the cancellation, or persistence, of infrared divergences is a formalism-dependent feature or a more robust property of the underlying theory.
} Therefore, building on the previous chapter, where we studied black-hole scattering in a higher-curvature, parity-violating deformation of GR, we now turn to a different class of beyond-GR effects: additional long-range fields and charges arising from a string-motivated low-energy effective theory. The central example in this chapter is Einstein--Maxwell--Dilaton (EMD) gravity, which provides a  setting in which to track how electromagnetic and dilatonic charges modify conservative scattering observables. As in the previous chapters, the broader motivation comes from precision gravitational-wave physics, where accurate gauge-invariant information about the conservative dynamics is needed to test departures from GR \cite{LIGOScientific:2016aoc,LIGOScientific:2016sjg,LIGOScientific:2016vlm,LIGOScientific:2017bnn,LIGOScientific:2019hgc,Purrer:2019jcp}.

From the effective-theory viewpoint, the origin of EMD is especially attractive because its extra fields are not introduced ad hoc. In the low-energy limit of heterotic string theory, the massless sector necessarily contains the graviton \(g_{\mu\nu}\), dilaton \(\theta\), antisymmetric tensor \(B_{\mu\nu}\), and non-Abelian gauge fields \(A_\mu^{I}\). For backgrounds that are independent of \textcolor{black}{\(d\) of the spacetime directions} and for gauge configurations commuting with \(p\) Abelian generators, the space of classical solutions enjoys a continuous \(O(d)\times O(d+p)\) symmetry that mixes \(g_{\mu\nu}\), \(B_{\mu\nu}\), \(\theta\), and \(A_\mu^{I}\) and acts as a solution-generating transformation \cite{Hassan:1991mq}. Acting with this “twist,” one can, for example, start from a ten-dimensional black 6-brane carrying magnetic charge but no electric or antisymmetric-tensor gauge charge and generate a family of inequivalent solutions with independent electric, magnetic, and antisymmetric-tensor charges \cite{Hassan:1991mq}. In this framework, additional degrees of freedom and higher-derivative corrections arise in tandem rather than as ad hoc additions, yielding symmetry-guided targets for EFT analyses and phenomenology. After restricting to the \(U(1)\) subsector, truncating the antisymmetric tensor, and transforming to Einstein frame, one arrives at the EMD theory used in this chapter; the resulting low-energy effective theory admits exact charged (and spinning) black-hole solutions \cite{GIBBONS1988741,Garfinkle:1990qj,Sen:1992ua}, making it a natural arena for probing possible stringy imprints in two-body dynamics.

For the hyperbolic and high-velocity regime, the post-Minkowskian description is the natural language, and scattering-amplitude methods provide an efficient way to isolate the classical conservative information \cite{Damour:2016gwp,Bini:2017wfr,Bini:2017xzy,Damour:2017zjx,Damour:2019lcq,Bini:2020flp,Bini:2020uiq,Damour:2020tta,Bini:2020rzn,Bini:2021gat,Bini:2022enm,Damour:2022ybd,Bini:2022wrq,Rettegno:2023ghr,Bini:2023fiz,Brandhuber:2022qbk,Brandhuber:2023hhy,Kosower:2018adc,DeAngelis:2023lvf,Cristofoli:2021vyo}. In this chapter we therefore adopt the scattering-amplitude approach to compute the conservative two-body potential, the eikonal phase, and consequently the scattering angle through one loop. Following \cite{Parra-Martinez:2020dzs}, we organize the calculation in terms of soft amplitudes and combine this with the Lippmann--Schwinger analysis of the potential, including the required EFT/Born subtraction of long-range iterations so that the final momentum-space potential is infrared finite. The chapter is organized as follows. In section~\eqref{ch4:sec2}, we motivate Einstein--Maxwell--Dilaton (EMD) theory as the low-energy effective description of the heterotic string, fix conventions and gauge choices, and derive the propagators and minimal set of on-shell Feynman rules required for $2\!\to\!2$ scattering. In section~\eqref{ch4:sec3}, we develop the analytic toolkit for the soft regime: using dimensional regularization, expansion by regions, and integration-by-parts (IBP) reduction, we obtain a compact basis of master integrals and their soft-limit expansions that capture the classical contributions. With these ingredients, we assemble the complete one-loop amplitudes across the relevant topologies (single exchanges, triangles, boxes, penguins) in the classical limit. In section~\eqref{ch4:sec4}, we extract the conservative two-body potential via the Lippmann-Schwinger equation and demonstrate that, after a careful EFT/Born subtraction of long-range iterations, the momentum-space potential is infrared finite. Section~\eqref{ch4:sec5} then determines the eikonal phase by exponentiating the momentum-space amplitude and, consequently, the scattering angle through one loop; we verify the smooth GR limit and delineate the separate roles of electromagnetic and dilatonic charges.




\section{Setup and methodology}\label{ch4:sec2}
\paragraph*{Low energy effective action of Heterotic string and emergence of EMD theory.}
The low-energy effective action of Heterotic string theory in string frame is given by \cite{Metsaev:1987zx,Hassan:1991mq},
\begin{align}
    \begin{split}
        S_{\textrm{low-eff}}=\int d^4x \sqrt{-G}\,e^{-\theta}\left[-R+\frac{1}{12}H_{\mu\nu\rho}H^{\mu\nu \rho}+G^{\mu\nu}\partial_{\mu}\theta \partial_{\nu}\theta-\frac{1}{8}F^2\right]
    \end{split}
\end{align}
    where $G_{\mu\nu}$ is the metric, $F_{\mu\nu}=\partial_\mu A_{\nu}-\partial_{\nu}A_{\mu}$ is the field strength corresponding to the Maxwell field $A_{\mu}$, and $\theta$ is the dilaton field. Moreover, $H_{\mu\nu\rho}=\partial_{\mu} B_{\nu\rho}+\textrm{cyclic per.}-(\Omega_3(A))_{\mu\nu\rho}$, where $B_{\mu\nu}$ is the antisymmetric tensor field and $(\Omega_3(A))_{\mu\nu\rho}=\frac{1}{4}(A_{\mu}F_{\nu\rho}+\textrm{cyclic per.})$ is the gauge Chern-Simons term. The theory has several noteworthy features. First, it arises from heterotic string theory compactified from ten to four dimensions. In the compactification procedure, the massless fields appearing in the effective action arise naturally. Also, only the $U(1)$ subsector of the full $E_{8}\times E_{8}$ or $\texttt{Spin}(32)/\mathbb{Z}_2$ theory has been retained. The metric $G_{\mu\nu}$, which naturally appears in the $\sigma$-model, is conformally related to the Einstein metric $g_{\mu\nu}$. Lastly, the action has been truncated at second order in derivatives of the fields. Now, ignoring the antisymmetric field $H_{\mu\nu\rho}$ and using an appropriate conformal field redefinition, the effective action takes the following form (restoring the factors of $m_p$),
    \begin{align}
        \begin{split}
              S_{\textrm{low-eff}}=\int d^4 x \sqrt{-g}\left[-\frac{m_p^2}{2}R+\frac{1}{2}g^{\mu\nu}\partial_{\mu}\theta\,\partial_\nu \theta-\frac{1}{16}\,e^{-\sqrt{2}\frac{\theta}{m_p}}F^2\right]\,.
        \end{split}
    \end{align}
The dilaton field $\theta$ plays a central role in the theory as its vacuum expectation value is related to the coupling of the theory as $\texttt{g}_{c}=e^{\sqrt{2}\langle\theta\rangle}$. The charged black hole solution (in unit of $G_N=1$) takes the following form \cite{Garfinkle:1990qj,GIBBONS1988741},
\begin{align}
\begin{split}
   & ds^2=\left(1-\frac{2M}{r}\right)dt^2-\left(1-\frac{2M}{r}\right)^{-1}dr^2-r\left(r-\frac{Q^2e^{-\sqrt{2}\theta_0}}{M}\right)d\Omega\,,
  \\ & e^{-\sqrt{2}\theta}=e^{-\sqrt{2}\theta_0}\left(1-\frac{Q^2\, e^{-\sqrt{2}\theta_0}}{Mr}\right)\,.
  \end{split}
\end{align}
Now, expanding the action around $\theta=\langle \theta\rangle:=\theta_0$ we  get,
     \begin{align}
        \begin{split}
              S_{\textrm{low-eff}}&=\int d^4 x \sqrt{-g}\left[-\frac{m_p^2}{2}R+\frac{1}{2}g^{\mu\nu}\partial_{\mu}\delta\theta\,\partial_\nu \delta\theta-\frac{1}{16}e^{-\sqrt{2}\theta_0/m_p}\,e^{-\sqrt{2}\frac{\delta\theta}{m_p}}F^2\right]\,, \\ &
              =\int d^4 x \sqrt{-g}\left[-\frac{m_p^2}{2}R+\frac{1}{2}g^{\mu\nu}\partial_{\mu}\delta\theta\,\partial_\nu \delta\theta-\frac{1}{16 \,\texttt{g}_{c}}\,e^{-\sqrt{2}\frac{\delta\theta}{m_p}}F^2\right]\,. \label{1.3h}
        \end{split}
    \end{align}
The theory described in \eqref{1.3h} has exact charged blackhole solution. The main goal of the paper is to study blackhole scattering and investigate how string theory modifies the fundamental classical observables such as the scattering angle (and the classical potential). As a standard procedure, we model the black holes by a charged scalar field $\Phi$ with dilaton dependent effective mass ($m_i(\theta)$)  , and the matter action is given by,
\begin{align}
    \begin{split}
        S_{\textrm{matter}}=\int d^4x \sqrt{-g}\sum_{i=1}^2\left[g^{\mu\nu}D_{\mu}\Phi^{\dagger}_iD_{\nu}\Phi_i-m_i^2(\theta)\Phi^{\dagger}_i\Phi_i\right]
    \end{split}
\end{align}
where the covariant derivative $D_{\mu}=\partial_{\mu}-ie\alpha_iA_{\mu}$ encodes the information about the gauge field $A_{\mu}$ and the electric charge of the scalar field $\Phi$, \textcolor{black}{ $Q_i=\alpha_i\,e$}.
Note that in the presence of the dilaton, the mass of the scalar field is now a function of the dilaton $\theta$. For computational purposes, the mass can be expanded around the background/asymptotic value of dilaton, i.e. $\theta_0$ as,
\begin{align}
    m_i(\theta)=m_i(\theta_0)+a_i(\theta_0)\,\delta\theta+b_i (\theta_0)\delta\theta^2+\cdots.
\end{align}
For notational simplicity, we use $\delta\theta\to \theta$, for further computation and define the path integral as, $$\int D\delta\theta\, D[A,g,\Phi_i]\cdots\to \int D\theta \,D[A,g,\Phi_i]\cdots\,. $$
\paragraph*{Propagators and Feynman rules.}
%\begin{align}
   % \begin{split}
      %  S_{\textrm{bulk}}=\int d^4 x \sqrt{-g}\left[-\frac{m_p^2}{2}R+\frac{1}{2}\partial_{\mu}\theta\,\partial^\mu \theta-\frac{\alpha'\,m_p^2}{16}\,e^{-\sqrt{2}\frac{\theta}{m_p}}F^2\right] 
   % \end{split}
%\end{align}
%and the matter action is given by,
%\begin{align}
   % \begin{split}
     %   S_{\textrm{matter}}=\int d^4x \sqrt{-g}\left[g^{\mu\nu}D_{\mu}\Phi^{\dagger}D_{\nu}\Phi-m^2(\theta)\Phi^{\dagger}\Phi\right]
   % \end{split}
%\end{align}
The two-point function for the Maxwell field will be modified due to the presence of string coupling $\texttt{g}_c$, but the two-point function of the dilaton and graviton remains unchanged. The gauge fixed action for the Maxwell field is given by,
\begin{align}
    \begin{split}
S_{\texttt{photon}}=\frac{1}{4\texttt{g}_c}\int d^4x \left(-\frac{1}{4}F^2-\frac{1}{2\xi}(\partial\cdot A)^2\right)\,.
    \end{split}
\end{align}
Immediately, we can identify the momentum space propagator in Feynman gauge as,
\begin{align}
\begin{split}
 \hspace{0cm}  D_{\mu\nu}(k)=\thesisinlinefeynman[0.075\linewidth]{photon_prop.png}\hspace{0.45cm}=
  4\texttt{g}_c\frac{-i\,\eta_{\mu\nu}}{k^2+i0}.
    \end{split}
\end{align}
Similarly the graviton\footnote{We define $P_{\mu\nu;\rho\sigma}=\left(\frac{1}{2}\eta_{\mu\rho}\eta_{\nu\sigma}+\frac{1}{2}\eta_{\mu\sigma}\eta_{\nu\rho}-\frac{1}{d-2}\eta_{\mu\nu}\eta_{\rho\sigma}\right)$.} and dilaton propagator takes the form,

\begin{align}
\begin{split}
  & 
 D_{\mu\nu;\rho \sigma}(k)= \thesisinlinefeynman[0.075\linewidth]{grav_propagator.png}\hspace{0.55cm}=\frac{i}{k^2+i0} \left(\frac{1}{2}\eta_{\mu\rho}\eta_{\nu\sigma}+\frac{1}{2}\eta_{\mu\sigma}\eta_{\nu\rho}-\frac{1}{d-2}\eta_{\mu\nu}\eta_{\rho\sigma}\right),\,\\ &
    D(k)=     \thesisinlinefeynman[0.075\linewidth]{dilaton_prop.png}\hspace{0.45cm}=\frac{i}{k^2+i0}.
  \end{split}
\end{align}

Now, expanding the bulk gravitational and matter actions about the background solution and truncating at cubic order yields the three-point interaction vertices needed for our analysis. We enumerate them below.
%\vspace{-2 cm}
{\scriptsize
\textbf{\textit{Bulk-Matter interaction vertex:}}
\vspace{0 cm}
\begin{align}
    \begin{split}
      &\hspace{0cm}     \bullet\,\,\thesisinlinefeynman[0.075\linewidth]{V2.png}\hspace{0.35cm}
\equiv 
\hspace{0 cm}\,\frac{i}{m_p}\left(k_3^{(\mu}k_2^{\nu)}-\frac{1}{2}\left(k_2\cdot k_3+m_0^2\right)\eta^{\mu\nu}\right),
%\\ &
 %\hspace{0cm}  
\bullet\,\, \thesisinlinefeynman[0.075\linewidth]{V3.png}\hspace{0.35cm}
\equiv 
\hspace{0 cm}\, i\,e\, \alpha_i\,(k_3^\mu-k_2^\mu),\\ &
 \hspace{0cm}     \,\,\bullet\,\,\thesisinlinefeynman[0.075\linewidth]{V5.png}\hspace{0.35cm}
\equiv 
\hspace{0 cm}\, -i \frac{2m_0 a}{m_p},\,\bullet\, \thesisinlinefeynman[0.075\linewidth]{4_1.png}\hspace{0.35cm}
\equiv 
\hspace{0 cm}\,2 {ie}^2 \eta ^{\mu \nu },\,\,\bullet\,\hspace{0cm} \thesisinlinefeynman[0.075\linewidth]{4_2.png}\hspace{0.35cm}
\equiv 
\hspace{0 cm}\,\frac{i e }{m_p}P^{\mu \nu ,\alpha \beta }(k_1-k_2)_{\beta },
    \end{split}
\end{align}
\begin{align}
    \begin{split}
    \\ &
\hspace{0cm}\bullet\, \thesisinlinefeynman[0.075\linewidth]{4_3.png}\hspace{0.35cm}
\equiv 
\hspace{0 cm}\,-\frac{i}{2 m_p^2}\Bigg[k_1^{\rho } k_2^{\sigma } \Big(-\eta _{\mu \nu } I_{\alpha \beta ,\rho \sigma }-\eta _{\alpha \beta } I_{\mu \nu ,\rho \sigma }+\eta _{\beta \nu } I_{\alpha \mu ,\rho \sigma }\\ &\hspace{0.8cm}+\eta _{\alpha \mu } I_{\beta \nu ,\rho \sigma }+\eta _{\beta \mu } I_{\alpha \nu ,\rho \sigma }+\eta _{\alpha \nu } I_{\beta \mu ,\rho \sigma }\Big)-P_{\mu \nu ,\alpha \beta }(k_1 \cdot k_2+m^2)\Bigg],\\ &
        \hspace{0cm}\bullet \thesisinlinefeynman[0.075\linewidth]{4_4.png}\hspace{0.35cm}
\equiv 
\hspace{0 cm}\,\frac{2 i \left(a^2-2 m_i\,b\right)}{m_p^2},\,\,\bullet \thesisinlinefeynman[0.075\linewidth]{4_5.png}\hspace{0.35cm}
\equiv 
\hspace{0 cm}\,-\frac{i\,a  m_i \eta ^{\mu \nu }}{m_p^2}.
    \end{split}
\end{align}
\textbf{\textit{Bulk  interaction vertex:}}
\begin{align}
    \begin{split}
      &\hspace{0 cm}     \bullet\,\,\thesisinlinefeynman[0.075\linewidth]{chichih.png}\hspace{0.35cm}
\equiv 
\hspace{0 cm}\,\frac{i}{4m_p}\left( k_1^\mu k_2^\nu + k_1^\nu k_2^\mu - k_1 \cdot k_2 ~\eta^{\mu \nu} \right)  ,
\\ &
\,\bullet\,\,\thesisinlinefeynman[0.075\linewidth]{AAchi.png}\hspace{0.35cm}
\equiv 
\hspace{0 cm}\, -i\frac{2\sqrt{2} }{16 \texttt{g}_c m_p}\left(\eta^{\mu\nu} k_2 \cdot k_3-k_2^{(\mu}k_3^{\nu)}\right),
\\ &
 \hspace{0 cm}     \bullet\,\,\thesisinlinefeynman[0.075\linewidth]{AAh.png}\hspace{0.35cm}
\equiv 
\hspace{0 cm}\,  -\frac{i  }{4\texttt{g}_c m_p}\Bigg[~  (k_2 \cdot k_3 ) P^{\alpha \beta , \mu \nu} +\eta^{\mu \nu} k_2^{ (\alpha} k_3^{\beta)}  
 + 
\frac{1}{2}\Big( \eta^{\alpha \beta}k_2^{\mu}k_3^{\nu} -\eta^{\alpha \mu}k_2^{\beta}k_3^{\nu} - \eta^{\alpha \nu}k_2^{\mu}k_3^{\beta}\\&\hspace{0.8cm} -  \eta^{\beta \mu}k_2^{\alpha}k_3^{\nu} - \eta^{\beta \nu}k_2^{\mu}k_3^{\alpha} \Big)    \Bigg].     
    \end{split}
\end{align}
}
With all the ingredients in hand, we now go on to the computation of the scattering amplitude at the soft limit and discuss how this particular limit of the full scattering amplitude is enough to generate the classically relevant observables. However, we also show that at one loop it can also generate \textit{some part} of the quantum correction to the classical potential. 
%\newpage
\section{Computation of scattering amplitude: Soft-loop expansion and classical limit}\label{ch4:sec3}
We first focus on the computation of the scattering amplitude in the soft limit \cite{Parra-Martinez:2020dzs} and then we go on to the computation of the Post-Minkowskian potential. A general $L$-loop, $n$-point amplitude takes the following form,
\begin{align}
\begin{split}
i \mathbfcal{A}_{n}^{L-\textrm{loop}}\left[\thesisampdiagram[0.15\linewidth]{npt.png}\right]=\sum_{j}\int \left(\prod_{l=1}^L\frac{d^d\ell_{l}}{(2\pi)^D}\right) \frac{1}{S_{j}}\frac{\mathcal{N}_{j}}{\prod_{\gamma_j}\rho_{\gamma_j}}
      \end{split}
\end{align}
where the sum runs over all possible L-loop Feynman diagrams $j$, for each diagram $\ell_{l}$ denotes the L loop momenta, $\gamma_j$ labels the propagators, and $S_j$ is associated with the symmetry factor of the diagram. $n_j's$ are the polynomials constructed from the irreducible scalar products constructed from external and loop momenta. In the subsequent computations, we will focus on the one-loop, four-point scattering amplitude, $\mathbfcal{A}_{4}^{1-\textrm{loop}}$, which is relevant for a binary scattering event. At the classical limit, any one-Loop, four-point amplitude can be written as,
\begin{align}
\begin{split}
   \mathbfcal{A}_{4}^{1-\textrm{loop}}={\mathfrak{c}_{\square}\mathcal{A}_{\square}+\mathfrak{c}_{\Join}\mathcal{A}_{\Join}+\mathfrak{c}_{\triangle}\mathcal{A}_{\triangle}+\mathfrak{c}_{\triangledown}\mathcal{A}_{\triangledown}+(\textrm{purely quantum})}\,.
   \end{split}
\end{align}
After integrating, it takes the form,
\begin{align}
     \mathbfcal{A}_{4}^{1-\textrm{loop}}=A+B \,q^2+\cdots+\frac{\alpha}{\sqrt{-q^2}}+\beta \log(\sqrt{-q^2})+\cdots\,.
\end{align}
Interestingly, it is trivial to show (by taking a suitable Fourier transform) that the non-analytic piece of the amplitude will contribute to the long-range interaction, which is relevant for classical physics. 
\\
\textcolor{black}{
\textbf{Soft expansion and the classical regime:}
Before turning to the detailed computations, we briefly review how the soft expansion emerges and how it encodes the classical limit of the amplitude \cite{Neill:2013wsa,Parra-Martinez:2020dzs}.
For hyperbolic scattering of two compact objects, the classical regime is characterized by an orbital angular momentum that is parametrically large compared to $\hbar$ (Bohr correspondence principle),
\begin{align}
J \gg \hbar \qquad (\text{equivalently, setting }\hbar=1:\; J\gg 1)\, .
\end{align}
Physically, $J\gg 1$ corresponds to a large impact parameter compared to the de Broglie wavelength, so that the scattering process admits a semiclassical description.
In this limit, the momentum transfer $q$ is small compared to the hard external scales set by the Mandelstam invariants and the masses, and classical contributions arise from the non-analytic (long-distance) dependence on $t=q^2$. More precisely, the large-$J$ limit induces a hierarchy of scales,
\begin{align}
s,\; |u|,\; m_i^2 \sim J^2 |t| \gg |t| = |q|^2 \, ,
\end{align}
so that expanding the amplitude at large $J$ is equivalent to expanding in small $|q|$.
This is what we refer to as the \emph{soft expansion}: it is an expansion around small momentum transfer, keeping the hard external kinematics fixed.
Within this expansion, the classical pieces of the amplitude can be systematically isolated (for instance, as the leading terms in $|q|$ at fixed $s$ and masses). It is useful to distinguish three classical limits, depending on how relativistic the scattering is:
\begin{align}
\begin{split}
    &\textrm{Generic classical limit:} \qquad\qquad\qquad\;\; J \gg 1,\\
    &\textrm{Near-static (PN) classical limit:} \qquad\,\,\, J \gg 1,\;\; \sigma \sim 1,\\
    &\textrm{High-energy classical limit:} \qquad\qquad\;\; J \gg 1,\;\; \sigma \gg 1,
\end{split}
\end{align}
where the boost parameter,
\begin{align}
\sigma \equiv \frac{k_1\!\cdot\! k_2}{m_1 m_2}
\end{align}
measures the relative velocity: $\sigma\simeq 1$ corresponds to the near-static/post-Newtonian regime, while $\sigma\gg 1$ captures ultra-relativistic scattering, where we have a stricter scale hierarchy: $s,|u|\gg m_i^2\sim J^2|t|\gg |t|=|q|^2$.
Since our goal is the conservative sector of the dynamics, we further restrict to the \emph{instantaneous} (or potential) long-range kinematics for the momentum transfer,
\begin{align}
q=(q^0,\boldsymbol{q})\,,\qquad |\boldsymbol{q}| \gg q^0\, .
\end{align}
This ensures that the long-range interaction is effectively instantaneous (no on-shell radiation is exchanged), which is precisely the regime relevant for defining a conservative potential from amplitudes. With these scalings in mind, loop integrals are conveniently analyzed using the method of regions.
For an internal graviton (or dilaton) momentum $\ell=(\omega,\boldsymbol{\ell})$, one encounters four standard momentum regions, distinguished by how $(\omega,\boldsymbol{\ell})$ compares to the hard scale $m$ and the soft scale $|\boldsymbol{q}|$:
\begin{align}
\begin{split}
    &\text{hard:}\qquad\qquad\;\; (\omega,\boldsymbol{\ell}) \sim (m,m),\\
    &\text{soft:}\qquad\qquad\;\; (\omega,\boldsymbol{\ell}) \sim (|\boldsymbol{q}|,|\boldsymbol{q}|) 
    \sim J^{-1}(m|\boldsymbol{v}|,\,m|\boldsymbol{v}|),\\
    &\text{potential:}\qquad\;\; (\omega,\boldsymbol{\ell}) \sim (|\boldsymbol{q}|\,|\boldsymbol{v}|,\,|\boldsymbol{q}|)
    \sim J^{-1}(m|\boldsymbol{v}|^2,\,m|\boldsymbol{v}|),\\
    &\text{radiation:}\qquad (\omega,\boldsymbol{\ell}) \sim (|\boldsymbol{q}|\,|\boldsymbol{v}|,\,|\boldsymbol{q}|\,|\boldsymbol{v}|)
    \sim J^{-1}(m|\boldsymbol{v}|^2,\,m|\boldsymbol{v}|^2).
\end{split}
\end{align}
Here we use the reference mass scale $m=m_1+m_2$ and the typical relative velocity scale $|\boldsymbol{v}|$ (which is related to $\sigma$).
The hard region captures short-distance physics and generates local (analytic) contributions, while the soft, potential, and radiation regions encode long-distance physics controlled by the small momentum transfer.
In particular, the \emph{potential} region is the one relevant for conservative dynamics: it corresponds to off-shell exchange with $\omega \ll |\boldsymbol{\ell}|$, matching the instantaneous regime $|\boldsymbol{q}|\gg q^0$.
The \emph{radiation} region, by contrast, is associated with near on-shell modes and is responsible for dissipative effects; we will therefore exclude it when focusing on the conservative potential. Except for the hard region, the remaining regions are governed by two small parameters: the soft scale $|\boldsymbol{q}|\sim J^{-1}$ controlling the classical expansion, and the velocity (equivalently $\sigma$) controlling the non-relativistic/PN/near-static expansion.
This double expansion will be the organizing principle for the computations that follow.}
\\\\
\textbf{Box topologies:}
To compute the amplitude from the box topologies, we encounter the following family of master integrals:
\begin{align}
    G^{\square}_{\alpha_1\alpha_2\alpha_3\alpha_4}=\int \frac{d^D\ell}{(2\pi)^D}\frac{1}{\rho_1^{\alpha_1}\rho_2^{\alpha_2}\rho_3^{\alpha_3}\rho_4^{\alpha_4}}
\end{align}
with, $\rho_1=\ell^2,\,\rho_2=(\ell-q)^2,\,\rho_3=(\ell+k_1)^2-m_1^2, \rho_4=(\ell-k_2)^2-m_2^2$. Now, for convenience, we choose the following kinematics,
\begin{align}
    \begin{split}
        k_1=\bar m_1 u_1-\frac{q}{2}, \,k_2=\bar m_2 u_2+\frac{q}{2}\implies k_1'=\bar m_1 u_1+\frac{q}{2},\,k_2'=\bar m_2 u_2-\frac{q}{2},\,
    \end{split}
\end{align}
Note that, by construction, the momentum transfer $q$ is orthogonal to the velocities $\bar u_is$. We would like to expand the full integral in the  soft region (which is important to take the classical limit), which follows the following hierarchy of scales,
\begin{align}
    |\ell|\sim |q|\ll m_i, \sqrt{s}\to (m_1+m_2)\,.
\end{align}
Now expanding the denominators ($D_2, \,D_4$) in this above mentioned limit we will get,
\begin{align}
\begin{split}
    &\rho_3=\ell^2+2 k_1\cdot \ell+i\epsilon= \ell^2+2 (\bar m_1u_1-\frac{q}{2})\cdot \ell+i\epsilon \sim 2 \bar m_1 u_1\cdot \ell+i\epsilon\,,\\ &
    \rho_4=\ell^2-2 k_2\cdot \ell+i\epsilon= \ell^2-2 (\bar m_2u_2+\frac{q}{2})\cdot \ell+i\epsilon \sim -2 \bar m_2 u_2\cdot \ell+i\epsilon\,.
    \end{split}
\end{align}
We parametrize the velocities as follows,
\begin{align}
    u_1=(1,0,0,0), \,u_2=(\sqrt{1+v^2},0,0,v).
\end{align}
Now in the soft region, the integral family reduces to,
\begin{align}
    \begin{split}
        \mathbfcal{G}^{}
_{\gamma_1\gamma_2\gamma_3\gamma_4}=\int \frac{d^d\ell}{(2\pi)^d}\frac{1}{D_1^{\gamma_1}D_2^{\gamma_2}D_3^{\gamma_3}D_4^{\gamma_4}}
    \end{split}
\end{align}
where the new ISPs (irreducible scalar products) are,
\begin{align}
    D_{i}\equiv \Big(2\bar m_1 u_1\cdot \ell+i\epsilon, -2 \bar m_2u_2\cdot \ell+i\epsilon,\ell^2,(\ell-q)^2\Big)\,.
\end{align}
After solving the IBP identities, one will get the following master integrals:
\begin{align}
    \vec f(x,\epsilon)=\{\mathbfcal{G}_{0,0,1,1},\mathbfcal{G}_{0,1,1,1},\mathbfcal{G}_{1,0,1,1},\mathbfcal{G}_{1,1,1,1}\}
\end{align}
where, $x=\sigma-\sqrt{\sigma^2-1}$ with, \textcolor{black}{$\sigma=\frac{k_1\cdot k_2}{m_1m_2}$}. We solve the master integrals by solving the following differential equation:
\begin{align}
    \partial_{x}\vec f(x,\epsilon)=A(x,\epsilon)\vec f(x,\epsilon)
\end{align}
where, the matrix $A(x,\epsilon)$ (using \cite{Gituliar:2017vzm}) is given by (only scale is $|q|$, so we set $q^2\to -1$, which can be restored from dimensional analysis),
\begin{align}
    \begin{split}
        A(x,\epsilon)=\left(
\begin{array}{ccc}
 0 & 0 & 0 \\
 0 & 0 & 0 \\
 \frac{2 (1-2 \epsilon )}{(x-1) (x+1)} & 0 & \frac{-x^2-1}{(x-1) x (x+1)} \\
\end{array}
\right)\,.
    \end{split}
\end{align}
The main goal is to find the canonical basis $ \vec g(x,\epsilon)=\mathbb{T}^{-1}(x,\epsilon)\vec f(x,\epsilon)$ such that the new differential equation will be $\epsilon$-factorized,   
\begin{align}
    \begin{split}
\partial_{x} \vec g(x,\epsilon)=\mathbb{T}^{-1}(A \mathbb{T}-\partial_x \mathbb{T})\vec g(x,\epsilon)\equiv \epsilon\, \mathbb{S}(x) \vec g(x,\epsilon)\,.
    \end{split}
\end{align}
The transformation matrix $\mathbb{T}$ and the $\epsilon$-factorized matrix $\mathbb{S}$ are given by,
\begin{align}
    \begin{split}
        \mathbb{T}(x,\epsilon)=\left(
\begin{array}{ccc}
 \frac{9 \epsilon }{2 \epsilon -1} & 0 & 0 \\
 2 & 1 & 0 \\
 \frac{2 x}{x^2-1} & -\frac{x}{x^2-1} & -\frac{3 x}{x^2-1} \\
\end{array}
\right),\,\,\,\mathbb{S}(x)=\left(
\begin{array}{ccc}
 0 & 0 & 0 \\
 0 & 0 & 0 \\
 \frac{6  }{x} & 0 & 0 \\
\end{array}
\right)\,.
    \end{split}
\end{align}
Therefore, the differential equation becomes,
\begin{align}
  &  d\,\vec{g}(x,\epsilon)=\epsilon\, \widetilde {\mathbb{S}}\,d\log(x)\,\vec g(x,\epsilon),\, \widetilde {\mathbb{S}}=\left(
\begin{array}{ccc}
 0 & 0 & 0 \\
 0 & 0 & 0 \\
 6 & 0 & 0 \\
\end{array}\right)\,.
\end{align}
  Finally, in  the Laporta basis, the formal solution takes the form,
\begin{center}
\includegraphics[width=0.42\linewidth]{basis.png}
\end{center}
Structurally, any one-loop soft-box contribution can be represented by the topology shown above.
To solve the master integrals, we need boundary conditions in the near-static ($\boldsymbol{v}\to 0$) limit.
Now we expand the soft-master integrals in the near static limit,
  \begin{align}
      \begin{split}
         \mathbfcal{G}^{\square}_{1,1,1,1}\xrightarrow[]{\boldsymbol{v}\to 0}\int_{-\infty}^{\infty} \frac{d\omega}{2\pi}\int \frac{d^{d-1}\ell}{(2\pi)^{d-1}} \frac{1}{ (2\bar m_1 \omega+i\epsilon)(-2\bar m_2 u_2^0\omega+2\bar m_2\boldsymbol{u_2}\cdot \boldsymbol{\ell}+i\epsilon)(\boldsymbol{-\ell^2}+i\epsilon)(\boldsymbol{-(\ell-q)^2}+i\epsilon)}\,.
      \end{split}
  \end{align}
Now integrating over $\omega$ (closing the contour in upper/lower half of $\omega$ plane) we get,
\begin{align}
    \begin{split}
       \hspace{0cm}\mathbfcal{G}^{\square}_{1,1,1,1}\xrightarrow[]{\boldsymbol{v}\to 0}&-\frac{i}{8\pi\bar m_1\bar m_2}\int \frac{d^d\boldsymbol{\ell}}{(2\pi)^{d-1}}\frac{1}{(\boldsymbol{u}_2\cdot \boldsymbol{\ell}+i\epsilon)(\boldsymbol{\ell}^2-i\epsilon)((\boldsymbol{\ell}-\boldsymbol{q})^2-i\epsilon)}\,,\\ &
        \to -\frac{1}{16\pi\bar m_1 \bar m_2}\int \frac{d^{d-1}\ell}{(2\pi)^{d-1}}\frac{\hat\delta(\boldsymbol{u}_2\cdot \boldsymbol{\ell})}{(\boldsymbol{\ell}^2-i\epsilon)((\boldsymbol{\ell-q})^2-i\epsilon)}\,,\\ &
        =-\frac{4^{\epsilon -2} \pi ^{\epsilon } \left(-q^2\right)^{-\epsilon -1} \Gamma (-\epsilon )^2 \Gamma (\epsilon +1)}{2\pi \bar{m}_1 \bar{m}_2 \Gamma (-2 \epsilon )}\,.
    \end{split}
\end{align}
Similarly, the near static limit of  $\mathbfcal{G}_{0,1,1,1}$ evaluates to,
\begin{align}
    \begin{split}
        \mathbfcal{G}^{\square}_{0,1,1,1}&\xrightarrow[]{v\to 0}-\frac{1}{4\pi\bar m_2}\int \frac{d^{d-1}\ell}{(2\pi)^{d-1}} \frac{1}{(\boldsymbol{\ell}^2-i\epsilon)[(\boldsymbol{\ell-\boldsymbol{q}})^2-i\epsilon]}\int_{-\infty}^{\infty}  \frac{d\omega}{\omega-\boldsymbol{u}_2\cdot \boldsymbol{\ell}-i\epsilon}\,,\\ &
        =-\frac{i\pi}{4\pi\bar m_2}\int 
    \frac{d^{d-1}\boldsymbol{\ell}}{(2\pi)^{d-1}}\frac{1}{(\boldsymbol{\ell}^2-i\epsilon)[(\boldsymbol{\ell-\boldsymbol{q}})^2-i\epsilon]}\,,\\ &
        =-\frac{i 16^{\epsilon -1} \pi ^{\epsilon +1} \left(-q^2\right)^{-\epsilon -\frac{1}{2}} \sec (\pi  \epsilon )}{2\pi\bar{m}_2 \Gamma (1-\epsilon )}
    \end{split}
\end{align}
and,
\begin{align}
    \begin{split}
        \mathbfcal{G}_{0,0,1,1}\xrightarrow[]{v\to 0}0\,.
    \end{split}
\end{align}
Now using (\ref{3.16}), the master integrals in the Laporta basis are given by,

\begin{center}
\begin{tcolorbox}[thesisresultbox, title=\textit{Master Integrals for box topologies}]
\restorethesisbodyformat

\begin{align}
    \begin{split}
     \hspace{0 cm}    &\mathbfcal{G}^{\square}_{0,0,1,1}(x,\epsilon)=0\,,\\&
        \mathbfcal{G}^{\square}_{0,1,1,1}(x,\epsilon)=- \textcolor{black}{\frac{1}{2 \pi}} \frac{i 16^{\epsilon -1} \pi ^{\epsilon +1} \sec (\pi  \epsilon )}{\bar{m}_2 \Gamma (1-\epsilon )}(-q^2)^{-\epsilon-\frac{1}{2}}\,,\\ &
        \mathbfcal{G}^{\square}_{1,0,1,1}(x,\epsilon)=- \textcolor{black}{\frac{1}{2 \pi}} \frac{i 16^{\epsilon -1} \pi ^{\epsilon +1} \sec (\pi  \epsilon )}{\bar{m}_1 \Gamma (1-\epsilon )}(-q^2)^{-\epsilon-\frac{1}{2}}\,,
        \\ &\mathbfcal{G}^{\square}_{1,1,1,1}(x,\epsilon)= \textcolor{black}{\frac{1}{2 \pi}} \frac{x\, 4^{2 \epsilon -1} \pi ^{\epsilon +\frac{3}{2}} \csc (\pi  \epsilon )}{\left(1-x^2\right) \bar{m}_1 \bar{m}_2 \Gamma \left(\frac{1}{2}-\epsilon \right)}(-q^2)^{-\epsilon-1}\,. 
    \end{split}
\end{align}
\end{tcolorbox}
\restorethesisbodyformat
\end{center}
\section*{Crossed-box topologies:}
In the crossed-box topologies, we encounter the following family of master integrals,
\begin{align}
    \begin{split}
        G^{\boxtimes}_{\gamma_1\gamma_2\gamma_3\gamma_4}=\int \frac{d^d\ell}{(2\pi)^d} \frac{1}{{\rho_1}^{\gamma_1}{\rho_2}^{\gamma_2}{\rho_4}^{\gamma_3}\tilde{\rho_4}^{\gamma_4}},
    \end{split}
\end{align}
with,  $\rho_1=\ell^2,\,\rho_2=(\ell-q)^2,\,\rho_3=(\ell+k_1)^2-m_1^2, \tilde\rho_4=(\ell+k_2-q)^2-m_2^2$. As mentioned earlier, we are interested in the soft/potential region of the amplitude. In the soft limit, the family reduces to,
\begin{align}
     \mathbfcal{G}^{\boxtimes}
_{\gamma_1\gamma_2\gamma_3\gamma_4}=\int \frac{d^d\ell}{(2\pi)^d}\frac{1}{\tilde D_1^{\gamma_1}\tilde D_2^{\gamma_2}\tilde D_3^{\gamma_3}\tilde D_4^{\gamma_4}}
\end{align}
where, 
$\tilde D_1=2\ell\cdot u_1+i\epsilon, \,\tilde D_2=2\ell\cdot u_2+i\epsilon,\, \tilde D_3=\ell^2,\,\tilde D_4=(\ell-q)^2. $ As we see in the crossed box topology, the two linearized matter propagators have the same sign, unlike the box integrals. As with the box topology, in this case also we have three master integrals: $\mathcal{G}_{0,0,1,1}^{\boxtimes},\,\mathcal{G}_{0,1,1,1}^{\boxtimes},\,\mathcal{G}_{1,1,1,1}^{\boxtimes}\,.$ Now, one will take a similar procedure as before to solve the integrals. However, one would notice that in the near static limit, the last will be zero.
\begin{align}\label{3.28} \mathbfcal{G}_{1,1,1,1}^{\boxtimes}=\frac{1}{4\bar{m}_1\bar{m}_2}\int \frac{d^{d-1}\boldsymbol{\ell}}{(2\pi)^{d-1}}\frac{1}{(\boldsymbol{\ell}^2-i\epsilon)\left[(\boldsymbol{\ell}-\boldsymbol{q})^2-i\epsilon\right]}\int \frac{d\omega}{2\pi} \frac{1}{(\omega+i\epsilon)(\omega+\boldsymbol{u}_2\cdot\boldsymbol{\ell}+i\epsilon)}\,.
\end{align}
The $\omega$ integral can be shown to be zero by closing the contour in the upper or lower half plane and computing the residues.
\begin{center}
\begin{tcolorbox}[thesisresultbox, title=\textit{Master Integrals for crossed-box topologies}]
\restorethesisbodyformat
\begin{centering}
\begin{align}
    \begin{split}
     \hspace{0 cm}    &\mathbfcal{G}^{\boxtimes}_{0,0,1,1}(x,\epsilon)=0\,,\\&
        \mathbfcal{G}^{\boxtimes}_{0,1,1,1}(x,\epsilon)=- \textcolor{black}{\frac{1}{2 \pi}} \frac{i 16^{\epsilon -1} \pi ^{\epsilon +1} \sec (\pi  \epsilon )}{\bar{m}_2 \Gamma (1-\epsilon )}(-q^2)^{-\epsilon-\frac{1}{2}}\,,\\ &
        \mathbfcal{G}^{\boxtimes}_{1,0,1,1}(x,\epsilon)=- \textcolor{black}{\frac{1}{2 \pi}} \frac{i 16^{\epsilon -1} \pi ^{\epsilon +1} \sec (\pi  \epsilon )}{\bar{m}_1 \Gamma (1-\epsilon )}(-q^2)^{-\epsilon-\frac{1}{2}}\,,\\&\mathbfcal{G}^{\boxtimes}_{1,1,1,1}(x,\epsilon)=0\,. 
    \end{split}
\end{align}
\end{centering}
\end{tcolorbox}
\restorethesisbodyformat
\end{center}
\noindent
Now we are ready to evaluate all the diagrams that contribute to the one-loop scattering amplitude in the soft limit. Below, we give the computational details. Before going to the one-loop computation we first give the results for tree level amplitude as it will play a crucial role to define IR finite potential.\\\\
%\newpage
\textbf{Tree-amplitude:}\\
\begin{center}
\includegraphics[width=0.36\linewidth]{Tree.png}
\end{center}
\begin{align}
\begin{split}
&i\mathbfcal{A}_{4}^{\texttt{tree}}
\\ &=\frac{1}{q^2}\left(-\frac{4  a^2 m_1 m_2}{m_p^2}+16  \alpha _1 \alpha _2 \texttt{g}_ce^2 m_1 m_2 \sigma  +\frac{ m_1^2 m_2^2 \left(-2 \sigma ^2 +1\right)}{2 m_p^2}\right)
    +\left(4  \alpha _1 \alpha _2 \texttt{g}_c e^2 -\frac{ 2 m_2 m_1 \sigma   }{4 m_p^2}\right)\,.    \end{split}
\end{align}
%\newpage
\noindent
We now move on to the computation of one-loop amplitudes in the soft limit and, consequently, use them to compute the classical Post-Minkowskian potential and the conservative scattering angle. \\\\
\textbf{One-loop amplitudes: }\\\\
In this section, we derive the one-loop amplitudes for different topologies and give the necessary details.  Expressions for the numerators of all the Feynman diagrams are given in the ancillary file \texttt{numerators\_emd.nb} submitted to the \texttt{arXiv} along with this paper. 
We start with computing the box (+ cross-box) amplitude.\\\\
\textbf{Soft amplitudes from Box (cross box) topologies:}\\
\textbullet $\,\,$ The amplitude for the box (+cross-box) diagram with double photon exchange is given by,
% \begin{align}
% \begin{split}
% &i\mathbfcal{A}_{4,\textrm{soft}}^{1-\textrm{loop}}\left[\begin{minipage}
%           [h]{0.40\linewidth}
% 	\vspace{1.2 pt}
% 	\scalebox{1}{\includegraphics[width=\linewidth]{Box_1.png}}
%       \end{minipage}\right]=\int_{\ell}\frac{\mathcal{N}^{(1)}_{\filledsquare{gray}}}{\rho_1 \rho_2 \rho_3 \rho_4}\Bigg|_{\textrm{soft}}\\ &
%      =\alpha _1^2 \alpha _2^2 e^4 2^{4 \epsilon -1} \pi ^{\epsilon } \texttt{g}_c^2 \left(-q^2\right)^{-\epsilon -\frac{3}{2}} \Bigg(8 \sigma  \left(\frac{4 \sqrt{\pi } m_1 m_2 \sqrt{-q^2} \sigma  \left(\sigma -\sqrt{\sigma ^2-1}\right) \csc (\pi  \epsilon )}{\left(-\sigma ^2+\sqrt{\sigma ^2-1} \sigma +1\right) \Gamma \left(\frac{1}{2}-\epsilon \right)}-\frac{i \left(m_1+m_2\right) q^2 \sec (\pi  \epsilon )}{\Gamma (1-\epsilon )}\right)\\ &+\frac{i \left(m_1+m_2\right) q^2 \sec (\pi  \epsilon ) \left(m_2^2 q^2 \sigma -m_1 m_2 q^2 (\sigma -4)+m_1^2 \sigma  \left(8 m_2^2+q^2\right)\right)}{m_1^2 m_2^2 \Gamma (1-\epsilon )}\Bigg]
%      \end{split}
% \end{align}

\begin{flalign}
& i\mathbfcal{A}_{4,\textrm{soft}}^{1-\textrm{loop}}\left[\thesisampdiagram[0.40\linewidth]{Box_1.png}\right]=\int_{\ell}\frac{\mathcal{N}^{(1)}_{\filledsquare{gray}}}{\rho_1 \rho_2 \rho_3 \rho_4}\Bigg|_{\textrm{soft}} +\int_{\ell}\frac{\mathcal{N}^{(1)}_{\Join}}{\rho_1 \rho_2 \rho_3 \tilde\rho_4}\Bigg|_{\textrm{soft}} \nonumber \\
      & = \frac{\alpha _1^2 \alpha _2^2 e^4 m_1 m_2 \sigma ^2 16^{\epsilon +1} \pi ^{\epsilon +\frac{1}{2}} \texttt{g}_c^2 \left(-q^2\right)^{-\epsilon -1} \csc (\pi  \epsilon )}{\sqrt{\sigma ^2-1} \Gamma \left(\frac{1}{2}-\epsilon \right)}\,. &
\end{flalign}
\\ 
\textbullet $\,\,$Amplitude corresponding to the box diagram having photon and graviton exchange is given by,      
%       \begin{align}
% \begin{split}
% &i\mathbfcal{A}_{4,\textrm{soft}}^{1-\textrm{loop}}\left[\begin{minipage}
%           [h]{0.40\linewidth}
% 	\vspace{1.2 pt}
% 	\scalebox{1}{\includegraphics[width=\linewidth]{Box_2.png}}
%       \end{minipage}\right]=\int_{\ell}\frac{\mathcal{N}^{(2)}_{\filledsquare{gray}}}{\rho_1 \rho_2 \rho_3 \rho_4}\Bigg|_{\textrm{soft}}\\ &
%       =     \frac{\alpha _1 \alpha _2 e^2 \texttt{g}_c}{128 m_1^3 m_2^3 \left(-q^2\right)^{3/2} m_p^2}\Bigg[m_1^4 m_2^4 16^{\epsilon +1} \pi ^{\epsilon } \left(-q^2\right)^{-\epsilon } \Big(-\frac{8 \sqrt{\pi } m_1 m_2 \sqrt{-q^2} \sigma  \left(2 \sigma ^2 (\epsilon -1)+1\right) \csc (\pi  \epsilon )}{\sqrt{\sigma ^2-1} (\epsilon -1) \Gamma \left(\frac{1}{2}-\epsilon \right)}\\ &+\frac{i \left(m_1+m_2\right) q^2 (2 \sigma  (3 \sigma  (\epsilon -1)+\epsilon )+1) \sec (\pi  \epsilon )}{(\epsilon -1) \Gamma (1-\epsilon )}\Big)\\ &+i \left(m_1+m_2\right) q^2 \Bigg(-m_2^2 q^6+m_1 m_2 q^6+2 m_1^3 m_2^3 q^2 \left(6 \sigma ^2-48 \sigma -1\right)\\ &-2 m_1^4 m_2^2 \left(6 \sigma ^2-1\right) \left(8 m_2^2+q^2\right)-m_1^2 \left(8 m_2^2 q^4+2 m_2^4 q^2 \left(6 \sigma ^2-1\right)+q^6\right)\Bigg) \Bigg]     
%       \end{split}
%       \end{align}

      \begin{align}
\begin{split}
&i\mathbfcal{A}_{4,\textrm{soft}}^{1-\textrm{loop}}\left[\thesisampdiagram[0.40\linewidth]{Box_2.png}\right]=\sum_{i=a,c}\int_{\ell}\frac{\mathcal{N}^{(2i)}_{\filledsquare{gray}}}{\rho_1 \rho_2 \rho_3 \rho_4}\Bigg|_{\textrm{soft}}+\sum_{j=b,d}\int_{\ell}\frac{\mathcal{N}^{(2j)}_{\Join}}{\rho_1 \rho_2 \rho_3 \tilde\rho_4}\Bigg|_{\textrm{soft}}\\ &
      =    \alpha _1 \alpha _2 e^2 m_1 m_2 2^{4 \epsilon -3} \pi ^{\epsilon } \texttt{g}_c \left(-q^2\right)^{-\epsilon } \left(\frac{-8 \sqrt{\pi } m_1 m_2 \sigma  \left(2 \sigma ^2 (\epsilon -1)+1\right) \csc (\pi  \epsilon )}{-q^2 \sqrt{\sigma ^2-1} (\epsilon -1) m_p^2 \Gamma \left(\frac{1}{2}-\epsilon \right)} \right.\\
      & \left. -\frac{i \left(m_1+m_2\right) \left(6 \sigma ^2 (\epsilon -1)+2 \sigma  \epsilon +1\right) \sec (\pi  \epsilon )}{\sqrt{-q^2} (\epsilon -1) m_p^2 \Gamma (1-\epsilon )}-\frac{i \left(m_1+m_2\right) \left(6 \sigma ^2 (\epsilon -1)-2 \sigma  \epsilon +1\right) \sec (\pi  \epsilon )}{\sqrt{-q^2} m_p^2 \Gamma (2-\epsilon )}\right).
      \end{split}
      \end{align}
Note that, a and b denotes the diagrams of the first row diagram in the first row and c and d denotes the diagrams of the second row.\\\\       
\textbullet $\,\,$The amplitude corresponding to the  box (+ cross-box) diagram with dilaton and photon exchange is given by,
%  \begin{align}
% \begin{split}
% &\hspace{0cm}i\mathbfcal{A}_{4,\textrm{soft}}^{1-\textrm{loop}}\left[\begin{minipage}
%           [h]{0.40\linewidth}
% 	\vspace{1.2 pt}
% 	\scalebox{1}{\includegraphics[width=\linewidth]{Box_3.png}}
%       \end{minipage}\right] =\int_{\ell}\frac{\mathcal{N}^{(3)}_{\filledsquare{gray}}}{\rho_1 \rho_2 \rho_3 \rho_4}\Bigg|_{\textrm{soft}}\\& = -\frac{a^2 \alpha _1 \alpha _2 e^2 2^{4 \epsilon -3} \pi ^{\epsilon } \texttt{g}_c \left(-q^2\right)^{-\epsilon -\frac{3}{2}}}{m_1^2 m_2^2 \left(-\sigma ^2+\sqrt{\sigma ^2-1} \sigma +1\right) m_p^2 \Gamma \left(\frac{1}{2}-\epsilon \right) \Gamma (1-\epsilon )}\Bigg[i m_2^3 q^4 \left(-\sigma ^2+\sqrt{\sigma ^2-1} \sigma +1\right) \sec (\pi  \epsilon ) \Gamma \left(\frac{1}{2}-\epsilon \right)\\ &+m_1^3 \left(-64 \sqrt{\pi } m_2^3 \sqrt{-q^2} \sigma  \left(\sqrt{\sigma ^2-1}-\sigma \right) \csc (\pi  \epsilon ) \Gamma (1-\epsilon )+i q^4 \left(-\sigma ^2+\sqrt{\sigma ^2-1} \sigma +1\right) \sec (\pi  \epsilon ) \Gamma \left(\frac{1}{2}-\epsilon \right)\right)\Bigg]
%       \end{split}
%       \end{align}

 \begin{flalign}
& i\mathbfcal{A}_{4,\textrm{soft}}^{1-\textrm{loop}}\left[\thesisampdiagram[0.40\linewidth]{Box_3.png}\right] =\sum_{i=a,c}\int_{\ell}\frac{\mathcal{N}^{(3i)}_{\filledsquare{gray}}}{\rho_1 \rho_2 \rho_3 \rho_4}\Bigg|_{\textrm{soft}}+\sum_{j=b,d}\int_{\ell}\frac{\mathcal{N}^{(3j)}_{\Join}}{\rho_1 \rho_2 \rho_3 \tilde\rho_4}\Bigg|_{\textrm{soft}} \nonumber \\
      & = -\frac{a^2 \alpha _1 \alpha _2 e^2 m_1 m_2 \sigma  2^{4 \epsilon +3} \pi ^{\epsilon +\frac{1}{2}} \texttt{g}_c \left(-q^2\right)^{-\epsilon -1} \csc (\pi  \epsilon )}{\sqrt{\sigma ^2-1} m_p^2 \Gamma \left(\frac{1}{2}-\epsilon \right)} . &
      \end{flalign}
\textbullet $\,\,$ The amplitude of the box (+cross-box) diagram corresponding to the double dilaton exchange is given by,
\begin{align}
\begin{split}
&i\mathbfcal{A}_{4,\textrm{soft}}^{1-\textrm{loop}}\left[\thesisampdiagram[0.40\linewidth]{Box_4.png}\right] =\int_{\ell}\frac{\mathcal{N}^{(4)}_{\filledsquare{gray}}}{\rho_1 \rho_2 \rho_3 \rho_4}\Bigg|_{\textrm{soft}}+\int_{\ell}\frac{\mathcal{N}^{(4)}_{\Join}}{\rho_1 \rho_2 \rho_3 \tilde\rho_4}\Bigg|_{\textrm{soft}} \\ &= \frac{a^4 2^{4 \epsilon -3} \pi ^{\epsilon +\frac{1}{2}} \left(-q^2\right)^{-\epsilon -1} \csc (\pi  \epsilon ) \left(8 m_1^2 m_2^2 \left(\sigma ^2-1\right)-q^2 \left(2 m_2 m_1 \sigma +m_1^2+m_2^2\right)\right)}{m_1 m_2 \left(\sigma ^2-1\right)^{3/2} m_p^4 \Gamma \left(\frac{1}{2}-\epsilon \right)}\,.
      \end{split}
      \end{align}
      \textbullet $\,\,$ The amplitude of the box (+cross-box) diagram corresponding to the graviton-dilaton  exchange is given by,
% \begin{align}
% \begin{split}
% &\hspace{-0 
% cm}i\mathbfcal{A}_{4,\textrm{soft}}^{1-\textrm{loop}}\left[\begin{minipage}
%           [h]{0.40\linewidth}
% 	\vspace{1.2 pt}
% 	\scalebox{1}{\includegraphics[width=\linewidth]{Box_5.png}}
%     \end{minipage}\right]    =\int_{\ell}\frac{\mathcal{N}^{(5)}_{\filledsquare{gray}}}{\rho_1 \rho_2 \rho_3 \rho_4}\Bigg|_{\textrm{soft}}\\ &
%     =\frac{a^2 \left(-q^2\right)^{-\epsilon -\frac{3}{2}}}{128 m_1 m_2 (\epsilon -1) m_p^4}\Bigg[m_1^2 m_2^2 2^{4 \epsilon +3} \pi ^{\epsilon } \Bigg(-\frac{4 \sqrt{\pi } m_1 m_2 \sqrt{-q^2} \left(\sqrt{\sigma ^2-1}-\sigma \right) \left(2 \sigma ^2 (\epsilon -1)+1\right) \csc (\pi  \epsilon )}{\left(-\sigma ^2+\sqrt{\sigma ^2-1} \sigma +1\right) \Gamma \left(\frac{1}{2}-\epsilon \right)}\\ &-\frac{i \left(m_1+m_2\right) q^2 (-2 \sigma +2 \sigma  \epsilon +\epsilon ) \sec (\pi  \epsilon )}{\Gamma (1-\epsilon )}\Bigg)\\ &+\frac{1}{\Gamma (1-\epsilon )}\Bigg(i \left(m_1+m_2\right) q^2 (16 \pi )^{\epsilon } \sec (\pi  \epsilon ) \Big(2 m_2^2 q^2 \sigma  (\epsilon -1)+m_1 m_2 q^2 (2 \sigma -2 \sigma  \epsilon +7 \epsilon -8)\\ &+2 m_1^2 \left(m_2^2 (-8 \sigma +8 \sigma  \epsilon -4 \epsilon )+q^2 \sigma  (\epsilon -1)\right)\Big)\Bigg)\Bigg]    \end{split}
%     \end{align}
\begin{flalign}
& i\mathbfcal{A}_{4,\textrm{soft}}^{1-\textrm{loop}}\left[\thesisampdiagram[0.40\linewidth]{Box_5.png}\right]    =\int_{\ell}\frac{\mathcal{N}^{(5)}_{\filledsquare{gray}}}{\rho_1 \rho_2 \rho_3 \rho_4}\Bigg|_{\textrm{soft}}+\int_{\ell}\frac{\mathcal{N}^{(5)}_{\Join}}{\rho_1 \rho_2 \rho_3 \tilde\rho_4}\Bigg|_{\textrm{soft}}\nonumber\\ &
    =\frac{a^2 m_1^2 m_2^2 2^{4 \epsilon -2} \pi ^{\epsilon +\frac{1}{2}} \left(-q^2\right)^{-\epsilon -1} \left(-2 \sigma ^2+2 \sigma ^2 \epsilon +1\right) \csc (\pi  \epsilon )}{\sqrt{\sigma ^2-1} (\epsilon -1) m_p^4 \Gamma \left(\frac{1}{2}-\epsilon \right)}  \nonumber \\
    & \hspace{0 cm} -\frac{i a^2 m_1 m_2 \left(m_1+m_2\right) \sigma  4^{2 \epsilon -1} \pi ^{\epsilon } \left(-q^2\right)^{-\epsilon -\frac{1}{2}} \sec (\pi  \epsilon )}{m_p^4 \Gamma (1-\epsilon )}\,. &
\end{flalign}
%\newpage
\textbullet $\,\,$   The amplitude of the box (+cross-box) diagram corresponding to the graviton-graviton  exchange is given by,  
%  \begin{align}
% \begin{split}
% &\hspace{-0 
% cm}i\mathbfcal{A}_{4,\textrm{soft}}^{1-\textrm{loop}}\left[\begin{minipage}
%           [h]{0.40\linewidth}
% 	\vspace{1.2 pt}
% 	\scalebox{1}{\includegraphics[width=\linewidth]{Box_6.png}}
%     \end{minipage}\right]    =\int_{\ell}\frac{\mathcal{N}^{(6)}_{\filledsquare{gray}}}{\rho_1 \rho_2 \rho_3 \rho_4}\Bigg|_{\textrm{soft}}\\ &
%     =\frac{4^{2 \epsilon -5} \pi ^{\epsilon } \left(-q^2\right)^{-\epsilon }}{(\epsilon -1)^2 m_p^4}\Bigg[-\frac{4 i \sec (\pi  \epsilon )}{\sqrt{-q^2} \Gamma (1-\epsilon )}\Bigg(4 m_1 m_2 \left(m_1+m_2\right) q^2 \sigma  (1-\epsilon ) (2 \sigma -2 \sigma  \epsilon +\epsilon )\\ &-\left(m_1+m_2\right) (2 \sigma -2 \sigma  \epsilon +\epsilon ) \left(2 m_1^2 m_2^2 \left(2 \sigma ^2 (\epsilon -1)+1\right)+q^4 (\epsilon -1)\right)\\ &+\frac{q^2 \left(2 m_1^3 \sigma  (\epsilon -1)+2 m_2^3 \sigma  (\epsilon -1)+m_2 m_1^2 (7 \epsilon -8)+m_2^2 m_1 (7 \epsilon -8)\right) \left(2 m_1^2 m_2^2 \left(2 \sigma ^2 (\epsilon -1)+1\right)+q^4 (\epsilon -1)\right)}{8 m_1^2 m_2^2}\Bigg)\\ &+(2 \sigma ^2 (\epsilon -1)+1)\times\\ &\Bigg(\frac{2 \sqrt{\pi } m_2 m_1 \left(\sqrt{\sigma ^2-1}-\sigma \right) \left(2 \sigma ^2 (\epsilon -1)+1\right) \csc (\pi  \epsilon ) \left(8 m_1^2 m_2^2 \left(\sigma ^2-1\right)-q^2 \left(2 m_2 m_1 \sigma +m_1^2+m_2^2\right)\right)}{q^2 \left(\sigma ^2-1\right) \left(\sigma  \left(\sqrt{\sigma ^2-1}-\sigma \right)+1\right) \Gamma \left(\frac{1}{2}-\epsilon \right)}\\ &+\frac{i m_1^2 \left(8 m_2^2+q^2\right) \sec (\pi  \epsilon ) \left(2 m_1 \sigma  (\epsilon -1)+m_2 \epsilon \right)}{\sqrt{-q^2} \Gamma (1-\epsilon )}+\frac{i m_2^2 \left(8 m_1^2+q^2\right) \sec (\pi  \epsilon ) \left(2 m_2 \sigma  (\epsilon -1)+m_1 \epsilon \right)}{\sqrt{-q^2} \Gamma (1-\epsilon )}\Bigg)\Bigg]    \end{split}
%     \end{align}
\begin{flalign}
&\hspace{0cm}i\mathbfcal{A}_{4,\textrm{soft}}^{1-\textrm{loop}}\left[\thesisampdiagram[0.40\linewidth]{Box_6.png}\right]    =\int_{\ell}\frac{\mathcal{N}^{(6)}_{\filledsquare{gray}}}{\rho_1 \rho_2 \rho_3 \rho_4}\Bigg|_{\textrm{soft}}+\int_{\ell}\frac{\mathcal{N}^{(6)}_{\Join}}{\rho_1 \rho_2 \rho_3 \tilde\rho_4}\Bigg|_{\textrm{soft}} \nonumber \\ & =\frac{m_1 m_2 2^{4 \epsilon -9} \pi ^{\epsilon +\frac{1}{2}} \left(-q^2\right)^{-\epsilon }  \csc (\pi  \epsilon ) \left(8 m_1^2 m_2^2 \left(\sigma ^2-1\right)-q^2 \left(2 m_2 m_1 \sigma +m_1^2+m_2^2\right)\right)}{-q^2 \left(2 \sigma ^2 (\epsilon -1)+1\right)^{-2} \left(\sigma ^2-1\right)^{3/2} (\epsilon -1)^2 m_p^4 \Gamma \left(\frac{1}{2}-\epsilon \right)} \nonumber \\ 
    & \hspace{1cm}+\frac{i m_1^2 m_2^2 \left(m_1+m_2\right) 4^{2 \epsilon -3} \pi ^{\epsilon } \epsilon  \left(-q^2\right)^{-\epsilon -\frac{1}{2}} \left(2 \sigma ^2 (\epsilon -1)+1\right) \sec (\pi  \epsilon )}{(\epsilon -1)^2 m_p^4 \Gamma (1-\epsilon )}\,.  &  
\end{flalign}
As one can immediately see by expanding the box(+cross-box) amplitudes, the $\frac{1}{\sqrt{-q^2}}$ terms cancelled, so they do not contribute to the classical potential.
\subsection*{Soft amplitudes from triangle topologies:}   
\textbullet  $\,\,$The soft amplitude corresponding to the triangle topology with photon lines is given by,

% \begin{align}
% \begin{split}
% &\hspace{0cm}i\mathbfcal{A}_{4,\textrm{soft}}^{1-\textrm{loop}}\left[\begin{minipage}
%           [h]{0.40\linewidth}
% 	\vspace{1.2 pt}
% 	\scalebox{0.8}{\includegraphics[width=\linewidth]{triangle_1.png}}
%     \end{minipage}\hspace{0cm}\right]=\int_{\ell}\frac{\mathcal{N}^{(1a)}_{\triangle}\rho_4+\mathcal{N}^{(1b)}_{\triangle}\rho_3}{\rho_1 \rho_2 \rho_3 \rho_4}\Bigg|_{\textrm{soft}}   \\ &\hspace{0 cm} =\frac{i \alpha _1^2 \alpha _2^2 e^4 \left(m_1+m_2\right) 4^{2 \epsilon +1} \pi ^{\epsilon } g_s^2 \left(-q^2\right)^{-\epsilon -\frac{1}{2}} \sec (\pi  \epsilon )}{\Gamma (1-\epsilon )} \,. 
%       \end{split}
%       \end{align}

\begin{align}
\thesisamplhs{i\mathbfcal{A}_{4,\textrm{soft}}^{1-\textrm{loop}}\left[\thesisampdiagram[0.32\linewidth]{triangle_1.png}\right]}
  &= \int_{\ell}\frac{\mathcal{N}^{(1a)}_{\triangle}\rho_4+\mathcal{N}^{(1b)}_{\triangle}\rho_3}{\rho_1 \rho_2 \rho_3 \rho_4}\Bigg|_{\textrm{soft}} \notag\\
\thesisamplhs{}
  &= \frac{i \alpha _1^2 \alpha _2^2 e^4 \left(m_1+m_2\right) 4^{2 \epsilon +1} \pi ^{\epsilon } g_s^2 \left(-q^2\right)^{-\epsilon -\frac{1}{2}} \sec (\pi  \epsilon )}{\Gamma (1-\epsilon )}\,.
\end{align}

\textbullet $\,\,$ The soft amplitude corresponding to the triangle topology with one graviton and one photon line is given by,   
\begin{align}
\thesisamplhs{i\mathbfcal{A}_{4,\textrm{soft}}^{1-\textrm{loop}}\left[\thesisampdiagram[0.32\linewidth]{triangle_4.png}\right]}
  &= \int_{\ell}\frac{\mathcal{N}^{(2a)}_{\triangle}}{\rho_1 \rho_2 \rho_3 } \Bigg|_{\textrm{soft}} \notag\\
\thesisamplhs{}
  &= -\frac{i \alpha _1 \alpha _2 e^2 2^{4 \epsilon -3} \pi ^{\epsilon } g_s \left(-q^2\right)^{-\epsilon -\frac{1}{2}} \sec (\pi  \epsilon )}{(\epsilon -1)^2 m_p^2 \Gamma (1-\epsilon )} \notag\\
\thesisamplhs{}
  &\phantom{=}\times \Big[m_1 q^2 \left(\epsilon ^2-\epsilon +1\right)+m_2 q^2 \sigma  \epsilon \notag\\
\thesisamplhs{}
  &\phantom{=}\qquad\;\; +8 m_1^2 m_2 \sigma  \left(\epsilon ^2-\epsilon +1\right)\Big].
\end{align}
and,
\begin{align}
\thesisamplhs{i\mathbfcal{A}_{4,\textrm{soft}}^{1-\textrm{loop}}\left[\thesisampdiagram[0.32\linewidth]{triangle_3.png}\right]}
  &= \int_{\ell}\frac{\mathcal{N}^{(2b)}_{\triangle}}{\rho_1 \rho_2 \rho_4 } \Bigg|_{\textrm{soft}} \notag\\
\thesisamplhs{}
  &= -\frac{i \alpha _1 \alpha _2 e^2 2^{4 \epsilon -3} \pi ^{\epsilon } g_s \left(-q^2\right)^{-\epsilon -\frac{1}{2}} \sec (\pi  \epsilon )}{(\epsilon -1)^2 m_p^2 \Gamma (1-\epsilon )} \notag\\
\thesisamplhs{}
  &\phantom{=}\times \Big[m_1 \sigma  \left(8 m_2^2 \left(\epsilon ^2-\epsilon +1\right)+q^2 \epsilon \right) \notag\\
\thesisamplhs{}
  &\phantom{=}\qquad\;\; +m_2 q^2 \left(\epsilon ^2-\epsilon +1\right)\Big]\,.
\end{align}
\textbullet $\,\,$ The soft amplitude corresponding to the triangle topology with graviton lines are given by,   
\begin{align}
\thesisamplhs{i\mathbfcal{A}_{4,\textrm{soft}}^{1-\textrm{loop}}\left[\thesisampdiagram[0.32\linewidth]{triangle_5.png}\right]}
  &= \int_{\ell}\frac{\mathcal{N}^{(3a)}_{\triangle}\rho_4+\mathcal{N}^{(3b)}_{\triangle}\rho_3}{\rho_1 \rho_2 \rho_3 } \Bigg|_{\textrm{soft}} \notag\\
\thesisamplhs{}
  &= \frac{i m_1^2 m_2^2 \left(m_1+m_2\right) 2^{4 \epsilon -6} \pi ^{\epsilon } \left(-q^2\right)^{-\epsilon -\frac{1}{2}} \sec (\pi  \epsilon )}{(\epsilon -1)^2 m_p^4 \Gamma (1-\epsilon )} \notag\\
\thesisamplhs{}
  &\phantom{=}\times \left(2 \sigma ^2 (\epsilon -1)^2-\epsilon \right)\,.
\end{align}
\textbullet $\,\,$ The soft amplitude corresponding to the triangle topology with dilaton lines is given by,   
\begin{align}
\thesisamplhs{i\mathbfcal{A}_{4,\textrm{soft}}^{1-\textrm{loop}}\left[\thesisampdiagram[0.32\linewidth]{triangle_6.png}\right]}
  &= \int_{\ell}\frac{\mathcal{N}^{(4a)}_{\triangle}\rho_4+\mathcal{N}^{(4b)}_{\triangle}\rho_3}{\rho_1 \rho_2 \rho_3 \rho_4 } \Bigg|_{\textrm{soft}} \notag\\
\thesisamplhs{}
  &= \frac{i a^2 2^{4 \epsilon -5} \pi ^{\epsilon } \left(-q^2\right)^{-\epsilon -\frac{1}{2}} \sec (\pi  \epsilon )}{m_1 m_2 m_p^4 \Gamma (1-\epsilon )} \notag\\
\thesisamplhs{}
  &\phantom{=}\times \left(8 a^2 m_1 m_2 \left(m_1+m_2\right)-32 b m_1^2 m_2^2\right)\,.
\end{align}
\textbullet $\,\,$ The soft amplitude corresponding to the triangle topology with one dilaton line and one graviton line is given by,   
% \begin{align}
% \begin{split}
% &\hspace{0cm}i\mathbfcal{A}_{4,\textrm{soft}}^{1-\textrm{loop}}\left[\begin{minipage}
%           [h]{0.50\linewidth}
% 	\vspace{1.2 pt}
% 	\scalebox{0.8}{\includegraphics[width=\linewidth]{triangle_8.png}}
%     \end{minipage}\hspace{0cm}\right]=\int_{\ell}\frac{\mathcal{N}^{(5a)}_{\triangle}}{\rho_1 \rho_2 \rho_3 } \Bigg|_{\textrm{soft}}    \\ &\hspace{0cm} =-\frac{i a^2 m_2 4^{2 \epsilon -3} \pi ^{\epsilon } (\epsilon +1) \left(8 m_1^2+q^2\right) \left(-q^2\right)^{-\epsilon -\frac{1}{2}} \sec (\pi  \epsilon )}{(\epsilon -1) m_p^4 \Gamma (1-\epsilon )},
%       \end{split}
%       \end{align} 

\begin{align}
\thesisamplhs{i\mathbfcal{A}_{4,\textrm{soft}}^{1-\textrm{loop}}\left[\thesisampdiagram[0.32\linewidth]{triangle_8.png}\right]}
  &= \int_{\ell}\frac{\mathcal{N}^{(5a)}_{\triangle}}{\rho_1 \rho_2 \rho_3 } \Bigg|_{\textrm{soft}} \notag\\
\thesisamplhs{}
  &= -\frac{i a^2 m_2 4^{2 \epsilon -3} \pi ^{\epsilon } (\epsilon +1) \left(8 m_1^2+q^2\right) \left(-q^2\right)^{-\epsilon -\frac{1}{2}} \sec (\pi  \epsilon )}{(\epsilon -1) m_p^4 \Gamma (1-\epsilon )}\,.
\end{align}

and,
% \begin{align}
% \begin{split}
% &\hspace{0cm}i\mathbfcal{A}_{4,\textrm{soft}}^{1-\textrm{loop}}\left[\begin{minipage}
%           [h]{0.50\linewidth}
% 	\vspace{1.2 pt}
% 	\scalebox{0.8}{\includegraphics[width=\linewidth]{triangle_7.png}}
%     \end{minipage}\hspace{0cm}\right] =\int_{\ell}\frac{\mathcal{N}^{(5b)}_{\triangle}}{\rho_1 \rho_2 \rho_3 } \Bigg|_{\textrm{soft}}   \\ &\hspace{0 cm} =-\frac{i a^2 m_1 4^{2 \epsilon -3} \pi ^{\epsilon } (\epsilon +1) \left(8 m_2^2+q^2\right) \left(-q^2\right)^{-\epsilon -\frac{1}{2}} \sec (\pi  \epsilon )}{(\epsilon -1) m_p^4 \Gamma (1-\epsilon )}.
%       \end{split}
%       \end{align}

\begin{align}
\thesisamplhs{i\mathbfcal{A}_{4,\textrm{soft}}^{1-\textrm{loop}}\left[\thesisampdiagram[0.32\linewidth]{triangle_7.png}\right]}
  &= \int_{\ell}\frac{\mathcal{N}^{(5b)}_{\triangle}}{\rho_1 \rho_2 \rho_3 } \Bigg|_{\textrm{soft}} \notag\\
\thesisamplhs{}
  &= -\frac{i a^2 m_1 4^{2 \epsilon -3} \pi ^{\epsilon } (\epsilon +1) \left(8 m_2^2+q^2\right) \left(-q^2\right)^{-\epsilon -\frac{1}{2}} \sec (\pi  \epsilon )}{(\epsilon -1) m_p^4 \Gamma (1-\epsilon )}\,.
\end{align}

\begingroup
\setlength{\thesisAmplitudeLhsWidth}{0.59\linewidth}
\setlength{\jot}{1.2pt}
\setlength{\abovedisplayskip}{2.5pt plus 1pt minus 1.5pt}
\setlength{\belowdisplayskip}{2.5pt plus 1pt minus 1.5pt}
\setlength{\abovedisplayshortskip}{1pt plus 1pt minus 0.5pt}
\setlength{\belowdisplayshortskip}{1pt plus 1pt minus 0.5pt}
\subsection*{Soft amplitudes from penguin topologies  }     
\textbullet $\,\,$ The soft amplitudes for penguin diagrams with two dilaton lines and one graviton line are given by,
\begin{align}
\thesisamplhs{i\mathbfcal{A}_{4,\textrm{soft}}^{1-\textrm{loop}}\left[\thesisampdiagram[0.42\linewidth]{penguin_1.png}\right]}
  &= \int_{\ell}\frac{\mathcal{N}_{\penguin}^{(1)}[\ell,q,\bar m_i,\sigma]}{q^2 \rho_1\rho_2\rho_3}\Bigg|_{\textrm{soft}} \notag\\
\thesisamplhs{}
  &= -\frac{i a^2 4^{2 \epsilon -5} \pi ^{\epsilon } \left(-q^2\right)^{-\epsilon -\frac{1}{2}} \sec (\pi  \epsilon )}{m_1 m_p^4 \Gamma (2-\epsilon )} \notag\\
\thesisamplhs{}
  &\phantom{=}\times \Big[m_2^2 q^2 \left(3 \sigma ^2-1\right) \notag\\
\thesisamplhs{}
  &\phantom{=}\qquad\;\; +2 m_1^2 \left(4 m_2^2 \left(\sigma ^2-1\right)+q^2\right) \notag\\
\thesisamplhs{}
  &\phantom{=}\qquad\;\; +4 m_1 m_2 q^2 \sigma\Big].
\end{align}
\vspace{-0.35\baselineskip}
\noindent and,
\begin{align}
\thesisamplhs{i\mathbfcal{A}_{4,\textrm{soft}}^{1-\textrm{loop}}\left[\thesisampdiagram[0.42\linewidth]{penguin_2.png}\right]}
  &= \int_{\ell}\frac{\mathcal{N}_{\penguin}^{(2)}[\ell,q,\bar m_i,\sigma]}{q^2 \rho_1\rho_2\rho_4}\Bigg|_{\textrm{soft}} \notag\\
\thesisamplhs{}
  &= -\frac{i a^2 4^{2 \epsilon -5} \pi ^{\epsilon } \left(-q^2\right)^{-\epsilon -\frac{3}{2}} \sec (\pi  \epsilon )}{m_2 m_p^4 \Gamma (2-\epsilon )} \notag\\
\thesisamplhs{}
  &\phantom{=}\times \Big[8 m_1 m_2^2 q^2 \left(m_1 \left(-\sigma ^2+2 \epsilon +1\right)-4 m_2 \epsilon \right) \notag\\
\thesisamplhs{}
  &\phantom{=}\qquad\;\; +q^4 \epsilon  \left(8 m_2^2+q^2\right)\Big]\,.
\end{align}
\vspace{-0.25\baselineskip}
\noindent\textbullet $\,\,$ The soft penguin amplitudes for with two photons, one graviton line are given by,
% \begin{align}
% \begin{split}
% &\hspace{0cm}i\mathbfcal{A}_{4,\textrm{soft}}^{1-\textrm{loop}}\left[\begin{minipage}
%           [h]{0.50\linewidth}
% 	\vspace{0.9 pt}
% 	\scalebox{0.8}{\includegraphics[width=\linewidth]{penguin_3.png}}    \end{minipage}\hspace{0cm}\right]=\int_{\ell}\frac{\mathcal{N}_{\penguin}^{(3)}[\ell,q,\bar m_i,\sigma]}{q^2 \rho_1\rho_2\rho_3}\Bigg|_{\textrm{soft}}    \\ &\hspace{0cm} =-\frac{i \alpha _1 e^2 4^{2 \epsilon -5} \pi ^{\epsilon } g_s \left(-q^2\right)^{-\epsilon -\frac{3}{2}} \sec (\pi  \epsilon )}{m_1^3 (\epsilon -1) m_p^2 \Gamma (2-\epsilon )}\Big[4 \left(2 \alpha _1+\alpha _2\right) m_1 m_2 q^6 \sigma  (\epsilon -1)^2+32 \left(2 \alpha _1+\alpha _2\right) m_1^3 m_2 q^4 \sigma  (\epsilon -1)^2\\ &-256 \alpha _2 m_1^5 m_2 q^2 \sigma  (\epsilon  (2 \epsilon -1)+1)+64 \alpha _1 m_1^4 \left(m_2^2 q^2 \left(\sigma ^2 (\epsilon -1) (4 \epsilon -3)+5 \epsilon -1\right)+q^4 \epsilon  (\epsilon +1)\right)\\ &-\alpha _1 q^2 (3 \epsilon -1) \left(2 m_2^2 q^4+q^6\right)+8 \alpha _1 m_1^2 \left(q^4 \left(m_2^2 (2-6 \epsilon )+q^2 \epsilon  (\epsilon +1)\right)-3 q^6 \epsilon +q^6\right)\Big]
%       \end{split}
%       \end{align}
\begin{align}
\thesisamplhs{i\mathbfcal{A}_{4,\textrm{soft}}^{1-\textrm{loop}}\left[\thesisampdiagram[0.42\linewidth]{penguin_3.png}\right]}
  &= \int_{\ell}\frac{\mathcal{N}_{\penguin}^{(3)}[\ell,q,\bar m_i,\sigma]}{q^2 \rho_1\rho_2\rho_3}\Bigg|_{\textrm{soft}} \notag\\
\thesisamplhs{}
  &= -\frac{i \alpha _1 e^2 m_1 m_2 16^{\epsilon -1} \pi ^{\epsilon } \texttt{g}_c \left(-q^2\right)^{-\epsilon -\frac{1}{2}} \sec (\pi  \epsilon ) }{(\epsilon -1) m_p^2 \Gamma (2-\epsilon )} \notag\\
\thesisamplhs{}
  &\phantom{=}\times \Big[ \alpha _1 m_2 \left(-3 \sigma ^2-4 \sigma ^2 \epsilon ^2\right) \notag\\
\thesisamplhs{}
  &\phantom{=}\qquad\;\; +\alpha _1 m_2 \left(\left(7 \sigma ^2-5\right) \epsilon +1\right) \notag\\
\thesisamplhs{}
  &\phantom{=}\qquad\;\; +4 \alpha _2 m_1 \sigma  \left(2 \epsilon ^2-\epsilon +1\right)\Big].
\end{align}
\vspace{-0.35\baselineskip}
\noindent and,
%\newpage
% \begin{align}
% \begin{split}
% &\hspace{0cm}i\mathbfcal{A}_{4,\textrm{soft}}^{1-\textrm{loop}}\left[\begin{minipage}
%           [h]{0.50\linewidth}
% 	\vspace{0.9 pt}
% 	\scalebox{0.8}{\includegraphics[width=\linewidth]{penguin_4.png}}    \end{minipage}\hspace{0cm}\right]=\int_{\ell}\frac{\mathcal{N}_{\penguin}^{(4)}[\ell,q,\bar m_i,\sigma]}{q^2 \rho_1\rho_2\rho_4}\Bigg|_{\textrm{soft}}    \\ &\hspace{0cm} =-\frac{i \alpha _2 e^2 \left(2 \sigma  \left(\sqrt{\sigma ^2-1}-\sigma \right)+1\right) 4^{2 \epsilon -5} \pi ^{\epsilon } g_s \left(-q^2\right)^{-\epsilon -\frac{3}{2}} \sec (\pi  \epsilon )}{m_2^3 \left(\sigma -\sqrt{\sigma ^2-1}\right)^2 (\epsilon -1)^2 m_p^2 \Gamma (1-\epsilon )} \Big[2 \alpha _2 m_1^2 \\ & \times\left(32 m_2^4 q^2 \left(\sigma ^2 (\epsilon -1) (4 \epsilon -3)+5 \epsilon -1\right)-q^4 (3 \epsilon -1) \left(8 m_2^2+q^2\right)\right)+4 m_2 m_1 \sigma  \Big(\left(\alpha _1+2 \alpha _2\right) q^4 (\epsilon -1)^2\\ &\times \left(8 m_2^2+q^2\right)-64 \alpha _1 m_2^4 q^2 (\epsilon  (2 \epsilon -1)+1)\Big)+\alpha _2 \left(8 m_2^2+q^2\right) \left(8 m_2^2 q^4 \epsilon  (\epsilon +1)+q^6 (1-3 \epsilon )\right)\Big]\,.
%       \end{split}
%       \end{align}
\begin{align}
\thesisamplhs{i\mathbfcal{A}_{4,\textrm{soft}}^{1-\textrm{loop}}\left[\thesisampdiagram[0.42\linewidth]{penguin_4.png}\right]}
  &= \int_{\ell}\frac{\mathcal{N}_{\penguin}^{(4)}[\ell,q,\bar m_i,\sigma]}{q^2 \rho_1\rho_2\rho_4}\Bigg|_{\textrm{soft}} \notag\\
\thesisamplhs{}
  &= -\frac{i \alpha _2 e^2 m_1 m_2 16^{\epsilon -1} \pi ^{\epsilon } \texttt{g}_c \left(-q^2\right)^{-\epsilon -\frac{1}{2}} \sec (\pi  \epsilon ) }{(\epsilon -1)^2 m_p^2 \Gamma (1-\epsilon )} \notag\\
\thesisamplhs{}
  &\phantom{=}\times \Big[ \alpha _2 m_1 \left(3 \sigma ^2+4 \sigma ^2 \epsilon ^2\right) \notag\\
\thesisamplhs{}
  &\phantom{=}\qquad\;\; +\alpha _2 m_1 \left(\left(5-7 \sigma ^2\right) \epsilon -1\right) \notag\\
\thesisamplhs{}
  &\phantom{=}\qquad\;\; -4 \alpha _1 m_2 \sigma  \left(2 \epsilon ^2-\epsilon +1\right)\Big]\,.
\end{align}
\vspace{-0.25\baselineskip}
\noindent\textbullet $\,\,$ The soft amplitudes for the penguin diagram with two photon lines and one dilaton line are given by,
% \begin{align}
% \begin{split}
% &\hspace{0cm}i\mathbfcal{A}_{4,\textrm{soft}}^{1-\textrm{loop}}\left[\begin{minipage}
%           [h]{0.50\linewidth}
% 	\vspace{0.9 pt}
% 	\scalebox{0.8}{\includegraphics[width=\linewidth]{penguin_5.png}}    \end{minipage}\hspace{0cm}\right]=\int_{\ell}\frac{\mathcal{N}_{\penguin}^{(5)}[\ell,q,\bar m_i,\sigma]}{q^2 \rho_1\rho_2\rho_3}\Bigg|_{\textrm{soft}}    \\ &\hspace{0cm} =\frac{i a \alpha _1 e^2 2^{4 \epsilon -\frac{17}{2}} \pi ^{\epsilon } g_s \left(-q^2\right)^{-\epsilon -\frac{3}{2}} \sec (\pi  \epsilon )}{m_1^3 m_p^2 \Gamma (1-\epsilon )}\\ 
%     & \hspace{0.8cm} \times \left[128 m_1^4 m_2 q^2 \left(\alpha _1-2 \alpha _2 \sigma \right)-q^4 \left(\alpha _1 m_2+2 \alpha _2 m_1\right) \left(8 m_1^2+q^2\right)\right]
%       \end{split}
%       \end{align}
\begin{align}
\thesisamplhs{i\mathbfcal{A}_{4,\textrm{soft}}^{1-\textrm{loop}}\left[\thesisampdiagram[0.42\linewidth]{penguin_5.png}\right]}
  &= \int_{\ell}\frac{\mathcal{N}_{\penguin}^{(5)}[\ell,q,\bar m_i,\sigma]}{q^2 \rho_1\rho_2\rho_3}\Bigg|_{\textrm{soft}} \notag\\
\thesisamplhs{}
  &= \frac{i a \alpha _1 e^2 2^{4 \epsilon -\frac{17}{2}} \pi ^{\epsilon } g_s \left(-q^2\right)^{-\epsilon -\frac{3}{2}} \sec (\pi  \epsilon )}{m_1^3 m_p^2 \Gamma (1-\epsilon )} \notag\\
\thesisamplhs{}
  &\phantom{=}\times \Big[128 m_1^4 m_2 q^2 \left(\alpha _1-2 \alpha _2 \sigma \right) \notag\\
\thesisamplhs{}
  &\phantom{=}\qquad\;\; -q^4 \left(\alpha _1 m_2+2 \alpha _2 m_1\right) \notag\\
\thesisamplhs{}
  &\phantom{=}\qquad\;\; \times \left(8 m_1^2+q^2\right)\Big].
\end{align}
\vspace{-0.35\baselineskip}
\noindent and,
% \begin{align}
% \begin{split}
% &\hspace{0cm}i\mathbfcal{A}_{4,\textrm{soft}}^{1-\textrm{loop}}\left[\begin{minipage}
%           [h]{0.50\linewidth}
% 	\vspace{0.9 pt}
% 	\scalebox{0.8}{\includegraphics[width=\linewidth]{penguin_6.png}}    \end{minipage}\hspace{0cm}\right]=\int_{\ell}\frac{\mathcal{N}_{\penguin}^{(6)}[\ell,q,\bar m_i,\sigma]}{q^2 \rho_1\rho_2\rho_4}\Bigg|_{\textrm{soft}}    \\ &\hspace{0cm} =-\frac{i a \alpha _2 e^2 2^{4 \epsilon -\frac{17}{2}} \pi ^{\epsilon } g_s \left(-q^2\right)^{-\epsilon -\frac{3}{2}} \sec (\pi  \epsilon )}{m_2^3 m_p^2 \Gamma (1-\epsilon )} \Big[128 m_1 m_2^4 q^2 \left(2 \alpha _1 \sigma -\alpha _2\right)+q^4 \left(2 \alpha _1 m_2+\alpha _2 m_1\right) \left(8 m_2^2+q^2\right)\Big]\,.
%       \end{split}
%       \end{align}
\begin{align}
\thesisamplhs{i\mathbfcal{A}_{4,\textrm{soft}}^{1-\textrm{loop}}\left[\thesisampdiagram[0.42\linewidth]{penguin_6.png}\right]}
  &= \int_{\ell}\frac{\mathcal{N}_{\penguin}^{(6)}[\ell,q,\bar m_i,\sigma]}{q^2 \rho_1\rho_2\rho_4}\Bigg|_{\textrm{soft}} \notag\\
\thesisamplhs{}
  &= -\frac{i a \alpha _2 e^2 2^{4 \epsilon -\frac{17}{2}} \pi ^{\epsilon } g_s \left(-q^2\right)^{-\epsilon -\frac{3}{2}} \sec (\pi  \epsilon )}{m_2^3 m_p^2 \Gamma (1-\epsilon )} \notag\\
\thesisamplhs{}
  &\phantom{=}\times \Big[128 m_1 m_2^4 q^2 \left(2 \alpha _1 \sigma -\alpha _2\right) \notag\\
\thesisamplhs{}
  &\phantom{=}\qquad\;\; +q^4 \left(2 \alpha _1 m_2+\alpha _2 m_1\right) \notag\\
\thesisamplhs{}
  &\phantom{=}\qquad\;\; \times \left(8 m_2^2+q^2\right)\Big]\,.
\end{align}
\vspace{-0.25\baselineskip}
\noindent\textbullet $\,\,$ The soft amplitudes for penguin topology with three graviton lines are given by,
% \begin{align}
% \begin{split}
% &\hspace{0cm}i\mathbfcal{A}_{4,\textrm{soft}}^{1-\textrm{loop}}\left[\begin{minipage}
%           [h]{0.50\linewidth}
% 	\vspace{1.2 pt}
% 	\scalebox{0.5}{\includegraphics[width=\linewidth]{penguin_7.png}}    \end{minipage}\hspace{0cm}\right] =\int_{\ell}\frac{\mathcal{N}_{\penguin}^{(7)}[\ell,q,\bar m_i,\sigma]}{q^2 \rho_1\rho_2\rho_3 \textcolor{black}{\rho_4}}\Bigg|_{\textrm{soft}}   \\ &\hspace{0cm} =\frac{i \left(m_1+m_2\right) 16^{\epsilon -3} \pi ^{\epsilon } \left(-q^2\right)^{-\epsilon -\frac{3}{2}} \sec (\pi  \epsilon )}{m_1 m_2 (\epsilon -1)^4 m_p^4 \Gamma (1-\epsilon )}\Big[8 m_1^2 m_2^2 q^4 (4 \sigma -2)+8 m_1^3 m_2 \left(m_2^2 q^2 \left(\sigma ^2+3\right)+2 q^4\right)\\ &+2 m_1 m_2 q^4 \left(8 m_2^2+2 q^2 \sigma +q^2\right)+\epsilon ^4 \Big(2 m_2^2 q^6+8 m_1^3 m_2 \left(m_2^2 q^2 \left(12 \sigma ^2-8\right)-q^4\right)+2 m_1^2 q^4 \left(4 m_2^2 (6 \sigma +3)+q^2\right)\\ &+m_1 m_2 q^4 \left(q^2 (6 \sigma +5)-8 m_2^2\right)\Big)+\epsilon ^3 \Big(-6 m_2^2 q^6+8 m_1^3 m_2 \left(m_2^2 q^2 \left(11-37 \sigma ^2\right)-4 q^4\right)+2 m_1^2 q^4 (4 m_2^2 (-22 \sigma -2)\\ &-3 q^2)+m_1 m_2 q^4 \left(-32 m_2^2-22 q^2 (\sigma +1)\right)\Big)+\epsilon ^2 \Big(6 m_2^2 q^6+8 m_1^3 m_2 \left(m_2^2 q^2 \left(39 \sigma ^2-15\right)+3 q^4\right)\\ &+m_1 m_2 q^4 \left(24 m_2^2+q^2 (30 \sigma +21)\right)+2 m_1^2 q^4 \left(4 m_2^2 (30 \sigma +3)+3 q^2\right)\Big)+\epsilon  \Big(-2 m_2^2 q^6+8 m_1^3 m_2 \Big(m_2^2 q^2 \left(1-15 \sigma ^2\right)\\ &-4 q^4\Big)+2 m_1^2 q^4 \left(4 m_2^2 (2-18 \sigma )-q^2\right)+m_1 m_2 q^4 \left(q^2 (-18 \sigma -10)-32 m_2^2\right)\Big)\Big]\,.
%       \end{split}
%       \end{align}
\begin{align}
\thesisamplhs{i\mathbfcal{A}_{4,\textrm{soft}}^{1-\textrm{loop}}\left[\thesisampdiagram[0.42\linewidth]{penguin_7.png}\right]}
  &= \int_{\ell}\frac{\mathcal{N}_{\penguin}^{(7a)}[\ell,q,\bar m_i,\sigma]}{q^2 \rho_1\rho_2\rho_3}\Bigg|_{\textrm{soft}} +\int_{\ell}\frac{\mathcal{N}_{\penguin}^{(7b)}[\ell,q,\bar m_i,\sigma]}{q^2 \rho_1\rho_2 \textcolor{black}{\rho_4}}\Bigg|_{\textrm{soft}} \notag\\
\thesisamplhs{}
  &= -\frac{i m_1^2 m_2^2 \left(m_1+m_2\right) 2^{4 \epsilon -9} \pi ^{\epsilon } \left(-q^2\right)^{-\epsilon -\frac{1}{2}}  \sec (\pi  \epsilon )}{(\epsilon -1)^4 m_p^4 \Gamma (1-\epsilon )} \notag\\
\thesisamplhs{}
  &\phantom{=}\times \Big[ -8 \epsilon ^4+11 \epsilon ^3-15 \epsilon ^2 \notag\\
\thesisamplhs{}
  &\phantom{=}\qquad\;\; +\sigma ^2 (\epsilon -1)^3 (12 \epsilon -1) \notag\\
\thesisamplhs{}
  &\phantom{=}\qquad\;\; +\epsilon +3\Big]\,.
\end{align}
\endgroup
\begin{tcolorbox}[thesisresultbox, title=\textit{Full one loop amplitude in small $|q|$ expansion}, boxsep=0.5mm, left=0.5mm, right=0.5mm, top=0.5mm, bottom=0.5mm]
\restorethesisbodyformat
\begingroup
\footnotesize
\setlength{\jot}{2pt}
\setlength{\abovedisplayskip}{4pt}
\setlength{\belowdisplayskip}{4pt}
\setlength{\abovedisplayshortskip}{2pt}
\setlength{\belowdisplayshortskip}{2pt}
\allowdisplaybreaks[4]
\begin{align*}
&\hspace{0cm}i\,\mathbfcal{A}_{\text{1-loop}}(\sigma,|q|,\ep)
= {} (-q^2)^{-(\tfrac32+\ep)}\;\\ & 
\hspace{0cm}\Biggl[
% ------------------------------------------------------------
% (0) Neutral / kinematic-only pieces (no a,e,\texttt{g}_c,\alpha_i)
% ------------------------------------------------------------
\;i(-q^2)\;
 \,\frac{2^{-9+4\ep}\,\pi^{\ep}\,\sec(\pi\ep)\,(m_1+m_2)}
             {(-1+\ep)^{3}\,\Gamma(1-\ep)\,\mpow{4}}
\Bigl(
      q^2\bigl(3-9\ep+2\ep^{2}-4\ep^{3}+4(-1+\ep)^{3}\sig-(-1+\ep)^{2}(-1+2\ep)\sig^{2}\bigr)
\\
&\hspace{0cm}
   m_1 m_2 +\,2 q^2\bigl(2 \epsilon ^3-\epsilon ^2+4 \epsilon -1\bigr)\,m_2^{2}
    +\,2 m_1^{2}\bigl(q^2(-1+\ep(4+\ep(-1+2\ep))) 
\\&\hspace{0cm} \,+ 4(-1+\ep)(-\ep+2(-1+\ep)^{2}\sig^{2})m_2^{2}\bigr)
\Bigr)+\; 
 i\,\frac{16^{-3+\ep}\,\pi^{\ep}\,\sec(\pi\ep)\,(m_1+m_2)}
             {(-1+\ep)^{4}\,\Gamma(1-\ep)\,m_1 m_2\,\mpow{4}}
\Bigl(
   2 q^{6}(-1+\ep)^{3}\ep\,m_2^{2} +\, q^{4} m_1 m_2
\\
&\hspace{0cm}
  \Bigl( 2 q^{2} + q^{2}\bigl(4 \sigma +(6 \sigma +5) \epsilon ^4-22 (\sigma +1) \epsilon ^3+3 (10 \sigma +7) \epsilon ^2-2 (9 \sigma +5) \epsilon\bigr)
   - 8\bigl(\epsilon ^4+4 \epsilon ^3-3 \epsilon ^2+4 \epsilon -2\bigr)m_2^{2}\Bigr)
\\
&\hspace{0cm}
 +\,2 q^{4} m_1^{2}\Bigl(q^{2}(-1+\ep)^{3}\ep
   + 4\bigl(-2+2\ep+3\ep^{2}-2\ep^{3}+3\ep^{4}+2(-1+\ep)^{3}(-2+3\ep)\sig\bigr)m_2^{2}\Bigr)
\\
&\hspace{0cm}
 +\,8 m_1^{3} m_2\Bigl(
     -q^{4}\bigl(\epsilon ^4+4 \epsilon ^3-3 \epsilon ^2+4 \epsilon -2\bigr)
     + q^{2}\bigl(3+\ep-15\ep^{2}+11\ep^{3}-8\ep^{4}+(-1+\ep)^{3}(-1+12\ep)\sig^{2}\bigr)m_2^{2}
   \Bigr)
\Bigr)
\\& \hspace{0cm}+\;
(-q^2)\;
\Bigl\{- i\,\frac{16^{-2+\ep}\,\pi^{\ep}\,\sec(\pi\ep)}
             {(-1+\ep)^{2}\,\Gamma(1-\ep)\,\mpow{4}}
\bigl(
  4 q^{2}(1-\ep)\sig(\ep+2\sig-2\ep\,\sig)\,m_1 m_2 (m_1+m_2)
\\
&\hspace{0cm}
 -\,(\ep+2\sig-2\ep\,\sig)(m_1+m_2)\bigl(q^{4}(-1+\ep) + 2(1+2(-1+\ep)\sig^{2})\,m_1^{2}m_2^{2}\bigr)
\\
&\hspace{0cm}
 +\,\frac{q^{2}\bigl(q^{4}(-1+\ep) + 2(1+2(-1+\ep)\sig^{2})\,m_1^{2}m_2^{2}\bigr)}
           {8\,m_1^{2}m_2^{2}}
   \bigl( 2(-1+\ep)\sig\,m_1^{3} + (-8+7\ep)\,m_1^{2} m_2
\\
&\hspace{0cm}
        +\,(-8+7\ep)\,m_1 m_2^{2} + 2(-1+\ep)\sig\,m_2^{3}\bigr)
\Bigr)\Big\}
\\
&\hspace{0cm} +\;
(-q^2)^{\tfrac32}\;
\frac{4^{-5+2\ep}\,\pi^{\ep}\,\bigl(1+2(-1+\ep)\sig^{2}\bigr)}
             {(-1+\ep)^{2}\,\mpow{4}}
\Biggl[
  \frac{i\,\sec(\pi\ep)\,(q^{2}+8 m_1^{2})\,m_2^{2}\,(\ep\,m_1+2(-1+\ep)\sig\,m_2)}{\sqrt{-q^2}\,\Gamma(1-\ep)}
\\
&\hspace{0cm}
 +\,\frac{i\,\sec(\pi\ep)\,m_1^{2}\,(2(-1+\ep)\sig\,m_1+\ep\,m_2)\,(q^{2}+8 m_2^{2})}{\sqrt{-q^2}\,\Gamma(1-\ep)}
\\
&\hspace{0cm}
 +\,\frac{2\sqrt{\pi}\,\bigl(1+2(-1+\ep)\sig^{2}\bigr)\,(-\sig+\SQ)\,\csc(\pi\ep)\,m_1 m_2\,
         \Bigl(8(\sig^{2}-1)m_1^{2}m_2^{2}-q^{2}(m_1^{2}+2\sig\,m_1 m_2 + m_2^{2})\Bigr)}
        {q^{2}\,(\sig^{2}-1)\,\bigl(1+\sig(-\sig+\SQ)\bigr)\,\Gamma(\tfrac12-\ep)}
\Biggr]
\\ &\hspace{0cm}\;+\; a^{4}\;
(-q^2)^{\tfrac12}\,
 \frac{2^{-3+4\ep}\,\pi^{\tfrac12+\ep}\,\csc(\pi\ep)\,
          \Bigl(8(\sig^2-1)\,m_1^{2}m_2^{2} - q^2\,(m_1^{2} + 2\sig\,m_1 m_2 + m_2^{2})\Bigr)}
          {(\sig^2-1)^{3/2}\,\Gamma(\tfrac12-\ep)\,m_1 m_2\,\mpow{4}}
% ------------------------------------------------------------
% (2) a^2 block
% ------------------------------------------------------------
\\ &\hspace{0cm}\;+\; a^{2}\Biggl[
 (-q^2)^{\tfrac32}\;
 \Bigl(-\,\frac{2^{-2+4\ep}\,\pi^{\tfrac12+\ep}\,\sqrt{-q^2}\,\bigl(1-2\sig^2+2\ep\,\sig^2\bigr)\,(-\sig+\SQ)\,\csc(\pi\ep)\,m_1^{2}m_2^{2}}
          {(-1+\ep)\,\bigl(1-\sig^2+\sig\,\SQ\bigr)\,\Gamma(\tfrac12-\ep)\,\mpow{4}}\Bigr)
\\
&\hspace{0cm}
 +\,(-q^2)\;
 \Bigl(-\, i\,\frac{2^{-4+4\ep}\,\pi^{\ep}\,q^{2}\,(\ep-2\sig+2\ep\,\sig)\,\sec(\pi\ep)\,m_1 m_2\,(m_1+m_2)}
             {(-1+\ep)\,\Gamma(1-\ep)\,\mpow{4}}\Bigr)
\\
&\hspace{0cm}
 +\,(-q^2)\;\Bigl[
 \begin{aligned}[t]
 &-\, i\,\frac{4^{-3+2\ep}\,\pi^{\ep}\,(1+\ep)\,\sec(\pi\ep)\,(q^2+8 m_1^{2})\,m_2}
              {(-1+\ep)\,\Gamma(1-\ep)\,\mpow{4}}
 \\
 &-\, i\,\frac{4^{-3+2\ep}\,\pi^{\ep}\,(1+\ep)\,\sec(\pi\ep)\,m_1\,(q^2+8 m_2^{2})}
              {(-1+\ep)\,\Gamma(1-\ep)\,\mpow{4}}
 \end{aligned}
 \Bigr]
\end{align*}
\refstepcounter{equation}\label{3.54j}
\begin{align*}
&\hspace{0cm}
 +\, i\,\frac{4^{-5+2\ep}\,\pi^{\ep}\,\sec(\pi\ep)}
             {\Gamma(2-\ep)\,m_1\,\mpow{4}}
 \Bigl(
   -q^4\,\ep\,(q^2+8 m_1^{2})
   + 8 q^2 m_1^{2} m_2\,(4\ep\,m_1 + (-1-2\ep+\sig^2)\,m_2)
 \Bigr)
\\
&\hspace{0cm}
 +\, i\,\frac{4^{-5+2\ep}\,\pi^{\ep}\,\sec(\pi\ep)}
             {\Gamma(2-\ep)\,m_2\,\mpow{4}}
 \Bigl(
   8 q^2 m_1 m_2^{2}\bigl((1+2\ep-\sig^2)\,m_1 - 4\ep\,m_2\bigr)
   + q^4 \ep\,(q^2+8 m_2^{2})
 \Bigr)
\\
&\hspace{0cm}
 +\,(-q^2)\;
 \Bigl[-\, i\,\frac{2^{-7+4\ep}\,\pi^{\ep}\,\sec(\pi\ep)\,(m_1+m_2)}
             {(-1+\ep)\,\Gamma(1-\ep)\,m_1 m_2\,\mpow{4}}
   \Bigl(q^2(-8+7\ep+2\sig-2\ep\,\sig)\,m_1 m_2
\\
&\hspace{0cm}
   +\,2 q^2(-1+\ep)\,\sig\,m_2^{2}
   +\,2 m_1^{2}\bigl(q^2(-1+\ep)\,\sig + (-4\ep-8\sig+8\ep\,\sig)\,m_2^{2}\bigr)\Bigr)\Bigr]
\\
&\hspace{0cm}
 +\,(-q^2)\;
 \frac{i\,2^{-5+4\ep}\,\pi^{\ep}\,\sec(\pi\ep)}
             {\Gamma(1-\ep)\,m_1 m_2\,\mpow{4}}
\Bigl[
\begin{aligned}
& q^{2} m_2 (a^{2}-2 b m_2)
 + a^{2} m_1 (q^{2}+8 m_2^{2}) \\
& + m_1^{2}\,\bigl(-2 b q^{2} + 8 m_2(a^{2}-4 b m_2)\bigr)
\end{aligned}
\Bigr]
\Biggr]
% ------------------------------------------------------------
% (3) e^2 \texttt{g}_c \alpha_1 \alpha_2 
% ------------------------------------------------------------
\\& \hspace{0cm}\;+\; e^{2}\texttt{g}_c\,\alpha_1\alpha_2\Biggl[
 (-q^2)\;
 \Bigl(-\, i\,\frac{2^{-3+4\ep}\,\pi^{\ep}\,\sec(\pi\ep)}
             {(-1+\ep)^{2}\,\Gamma(1-\ep)\,\mpow{2}}\,
   \bigl( q^{2}(1-\ep+\ep^{2})\,m_1 + q^{2}\ep\,\sig\,m_2 + 8(1-\ep+\ep^{2})\,\sig\,m_1^{2} m_2 \bigr)\Bigr)
\\
&\hspace{0cm}
 +\,
 \frac{2^{-3+4\ep}\,\pi^{\ep}\,m_1 m_2}
             {(-1+\ep)\,\mpow{2}}
\Biggl[
  -\,\frac{8\sqrt{\pi}\,\sqrt{-q^2}\,\sig\bigl(1+2(-1+\ep)\sig^{2}\bigr)\,\csc(\pi\ep)\,m_1 m_2}{\SQ\,\Gamma(\tfrac12-\ep)}
  \\ &
  \hspace{0cm}+\,\frac{i\,q^{2}\bigl(1+2\sig(\ep+3(-1+\ep)\sig)\bigr)\,\sec(\pi\ep)\,(m_1+m_2)}{\Gamma(1-\ep)}
\Biggr]
\\
&\hspace{0cm}
 +\,(-q^2)\;
 \frac{i\,2^{-7+4\ep}\,\pi^{\ep}\,\sec(\pi\ep)\,(m_1+m_2)}
             {\Gamma(2-\ep)\,m_1^{3} m_2^{3}\,\mpow{2}}
\Bigl(
 16(1-2\ep\,\sig+6(-1+\ep)\sig^{2})\,m_1^{4} m_2^{4}
\\
&\hspace{0cm}
 +\,2 q^{2} m_1^{2} m_2^{2}\Bigl((1+6(-1+\ep)\sig^{2})\,m_1^{2}
    - \bigl(1-6(-8+\sig)\sig + \ep(8-46\sig+6\sig^{2})\bigr)\,m_1 m_2
    + (1+6(-1+\ep)\sig^{2})
\\
&\hspace{0cm}
\,m_2^{2}\Bigr) +\, q^{4}(-1+\ep)\,\Bigl(8 m_1^{2} m_2^{2} + q^{2}(m_1^{2}-m_1 m_2 + m_2^{2})\Bigr)
\Bigr)
\Biggr]
% ------------------------------------------------------------
% (4) a^2 e^2 \texttt{g}_c \alpha_1 \alpha_2 
% ------------------------------------------------------------
\\&\hspace{0cm}\;+\; a^{2} e^{2} \texttt{g}_c\,\alpha_1\alpha_2\Biggl[
 \frac{2^{3+4\ep}\,\pi^{\tfrac12+\ep}\,\sqrt{-q^2}\,\sig(\sig-\SQ)\,\csc(\pi\ep)\,m_1 m_2}
      {\bigl(-1+\sig^{2}-\sig\,\SQ\bigr)\,\Gamma(\tfrac12-\ep)\,\mpow{2}}
\\
&\hspace{0cm}
 +\,\frac{i\,2^{4\ep}\,\pi^{\ep}\,q^{2}\,\sec(\pi\ep)\,(m_1+m_2)}{\Gamma(1-\ep)\,\mpow{2}}
 \;+\;
 (-q^2)\;
 \frac{i\,2^{-3+4\ep}\,\pi^{\ep}\,\sec(\pi\ep)\,(m_1+m_2)}
             {\Gamma(1-\ep)\,m_1^{2} m_2^{2}\,\mpow{2}}
\Bigl(-q^{2} m_1 m_2 + q^{2} m_2^{2} 
% ------------------------------------------------------------
% (5) e^4 \texttt{g}_c^2 \alpha_1^2 \alpha_2^2 
% ------------------------------------------------------------
\\&\hspace{0cm}\;+ m_1^{2}(q^{2}+8 m_2^{2})\Bigr)
\Biggr]+\; e^{4} \texttt{g}_c^{2}\alpha_1^{2}\alpha_2^{2}\Biggl[
 (-q^2)\; i\,\frac{4^{1+2\ep}\,\pi^{\ep}\,\sec(\pi\ep)\,(m_1+m_2)}{\Gamma(1-\ep)}
\\
&\hspace{0cm}
 +\,4^{1+2\ep}\,\pi^{\ep}\,\sig
 \Biggl(
    \frac{4\sqrt{\pi}\,\sqrt{-q^2}\,\sig\,(\sig-\SQ)\,\csc(\pi\ep)\,m_1 m_2}{\bigl(1-\sig^{2}+\sig\,\SQ\bigr)\,\Gamma(\tfrac12-\ep)}
   -\frac{i\,q^{2}\,\sec(\pi\ep)\,(m_1+m_2)}{\Gamma(1-\ep)}
 \Biggr)
\\
&\hspace{0cm}
 +\,(-q^2)\;
 \Bigl[-\, i\,\frac{2^{-1+4\ep}\,\pi^{\ep}\,\sec(\pi\ep)\,(m_1+m_2)}{\Gamma(1-\ep)\,m_1^{2} m_2^{2}}\,
   \bigl(-q^{2}(-4+\sig)\,m_1 m_2 + q^{2}\sig\,m_2^{2} + \sig\,m_1^{2}(q^{2}+8 m_2^{2})\bigr)\Bigr]
\Biggr]
% ------------------------------------------------------------
% (6) linear in \alpha_1 or \alpha_2  (with e^2 \texttt{g}_c)
% ------------------------------------------------------------
\\&\hspace{0cm}\;+\; e^{2} \texttt{g}_c \Biggl\{
\alpha_1\,
\Bigl[
 \; i\,\frac{4^{-5+2\ep}\,\pi^{\ep}\,\sec(\pi\ep)}
             {\Gamma(2-\ep)\,m_1^{3} m_2^{3}\,\mpow{2}}
 \bigl(
  -q^{2}(-1+3\ep)\,(q^{6}+2 q^{4} m_2^{2})
 \\
 \displaybreak[3]
 &\hspace{0cm}
  + 64 m_1^{4}\Bigl(
      q^{4}\ep(1+\ep)
      + q^{2}(-1+5\ep + (-1+\ep)(-3+4\ep)\sig^{2})\,m_2^{2}
    \Bigr)
\\
 &\hspace{0cm}
  + 8 m_1^{2}\Bigl(
      q^{6}-3 q^{6}\ep
      + q^{4}(q^{2}\ep(1+\ep)+(2-6\ep)m_2^{2})
    \Bigr)
 \bigr)
\Bigr]
\\
& \hspace{0cm} +\;
\alpha_2\,
\Bigl[
 \frac{i\,4^{-5+2\ep}\,\pi^{\ep}\,\sec(\pi\ep)}
      {(-1+\ep)^{2}(\sig-\SQ)^{2}\,\Gamma(1-\ep)\,m_2^{3}\,\mpow{2}}\,
 \bigl(1+2\sig(-\sig+\SQ)\bigr)
\\
&\hspace{0cm}
 \times\Bigl(
   (q^{2}+8 m_2^{2})\bigl(q^{6}(1-3\ep)+8 q^{4}\ep(1+\ep)\,m_2^{2}\bigr)\,\alpha_2
 +\,2 m_1^{2}\bigl(32 q^{2}\,m_2^{4} - q^{4}(1-3\ep)(q^{2}+8 m_2^{2})\bigr)\,\alpha_2
\\
&\hspace{0cm}
 +\,4 \sig\,m_1 m_2\bigl(-64 q^{2}(1+\ep(-1+2\ep))\,m_2^{4}\,\alpha_1
        + q^{4}(-1+\ep)^{2}(q^{2}+8 m_2^{2})(\alpha_1+2\alpha_2)\bigr)
 \Bigr)
\Bigr]
\Biggr\}
% ------------------------------------------------------------
% (7) linear in a :  a e^2 \texttt{g}_c (\alpha_1,\alpha_2)
% ------------------------------------------------------------
\\ &\hspace{0cm}\;+\; a\,e^{2} \texttt{g}_c \Biggl[
 -\, \frac{i\,2^{-\tfrac{17}{2}+4\ep}\,\pi^{\ep}\,\sec(\pi\ep)}
             {\Gamma(1-\ep)\,m_2^{3}\,\mpow{2}}\,
 \alpha_2 \Bigl( 128 q^{2} m_1\,m_2^{4}\,(2\sig\,\alpha_1 - \alpha_2)
       + q^{4}(q^{2}+8 m_2^{2})\,(2 m_2 \alpha_1 + m_1 \alpha_2) \Bigr)
\\
&\hspace{0cm}
 +\, \frac{i\,2^{-\tfrac{17}{2}+4\ep}\,\pi^{\ep}\,\sec(\pi\ep)}
             {\Gamma(1-\ep)\,m_1^{3}\,\mpow{2}}\,
 \alpha_1 \Bigl(
   128 q^{2} m_1^{4} m_2\,(\alpha_1-2\sig\,\alpha_2)
\\
&\hspace{0cm}
       - q^{4}(q^{2}+8 m_1^{2})\,(m_2 \alpha_1 + 2 m_1 \alpha_2)
 \Bigr)
\Biggr]
    +\mathcal{O}(q^2 \log(-q^2))\Biggr]\,.\tag{\theequation}
\end{align*}
\endgroup
\end{tcolorbox}
\restorethesisbodyformat
\section{Post-Minkowskian potential: The Lippmann-Schwinger equation and the Infrared subtraction} \label{ch4:sec4}   
The Post-Minkowskian potential can be derived from the Lippmann-Schwinger equation. The computation will be easier in the center-of-mass coordinate. The kinematics we choose is the following,
\begin{align}
\begin{split}
&   \texttt{incoming momenta:}\,\, k_1^\mu=(E_1,\boldsymbol{k}),\,k_2^\mu=(E_2,-\boldsymbol{k}),\\ &
\texttt{outgoing momenta:} k_1^{'\mu}=(E_1,\boldsymbol{k+q}), \,k_2^{'\mu}=(E_2,\boldsymbol{-k-q})\,.
\end{split}
\end{align}
One can define the potential in terms of scattering amplitude using the Lippmann-Schwinger equation, which states,
\begin{align}
    \mathbfcal{A}(\boldsymbol{k},\boldsymbol{k'})=V(\boldsymbol{k},\boldsymbol{k'})+\int d^3\boldsymbol{\ell} \,\mathbfcal{A}(\boldsymbol{k},\boldsymbol{\ell}) \,G(\boldsymbol{k},\boldsymbol{\ell})\,V(\boldsymbol{\ell},\boldsymbol{k'}),\label{4.2f}
\end{align}
where $G(\boldsymbol{k},\boldsymbol{\ell})$ is the Green function, and the choice is somewhat arbitrary. However, consistency conditions on scattering amplitude, such as unitarity and demanding that the potential should be a real one, can show (using the optical theorem) that,
\begin{align}
\texttt{Im}\,G(\boldsymbol{k},\boldsymbol{\ell})=\frac{\delta(|\boldsymbol{\ell}|^2-|\boldsymbol{k}|^2)}{16\pi^2 \sqrt{s}}.
\end{align}
The simplest choice of the Green function that satisfies the above is the following 

\begin{align}
\begin{split}
G(\boldsymbol{k},\boldsymbol{\ell})=\frac{1}{(2\pi)^3}\frac{1}{2\sqrt{s}}\frac{1}{|\boldsymbol{\ell}|^2-|\boldsymbol{k}|^2-i0}+\cdots\,.\label{4.4g}
   \end{split}
\end{align}
\textcolor{black}{In principle, one may add to \eqref{4.4g} any term that is regular, provided the constraint is satisfied. Different choices correspond to alternative effective potentials $V$ that differ only by off-shell pieces \cite{Correia:2024jgr}. 
Importantly, the constant piece will lead to a modification of the effective potential in the classical regime. In our computation, we fix the constant piece by comparing it with the traditional Lippmann-Schwinger Green's function, and we will comment on that in detail in the subsequent discussions.}
%However, the the other higher order analytic pieces may lead to short range interaction (or they may vanish in the soft limit as they result scaleless integrals) which are irrelevant for classical long range force. We will explicitly illustrate this at the end of this section.  Nevertheless, at very large center of mass energy those contributions will be more and more sub-leading. In our computation we fix the constant piece by comparing it with the traditional Lippmann-Schwinger Green's function and we will comment on that in detail in the subsequent discussions.}
\begin{figure}
 \hspace{0cm}   \centering    \includegraphics[width=0.9\linewidth]{LS.png}
    \caption{Diagrammatic representation of Lippmann-Schwinger equation}
    \label{fig1}
\end{figure}
Now, by inverting \eqref{4.2f} and truncating at one-loop order, we can write the potential in momentum space iteratively in terms of the scattering amplitude (see Fig.~\ref{fig1} for the diagrammatic representation),
\begin{align}
    V(\boldsymbol k,\boldsymbol k')\Big|_{\texttt{1-loop}}=\mathbfcal{A}_{\texttt{1-loop}}(\boldsymbol k,\boldsymbol k')-\int d^3\boldsymbol{\ell}\,\mathbfcal{A}_{\texttt{tree}}(\boldsymbol{k},\boldsymbol{\ell}) G(\boldsymbol{k},\boldsymbol{\ell})\mathbfcal{A}_{\texttt{tree}}(\boldsymbol{\ell},\boldsymbol{k'})\,. \label{4.5h}
\end{align}
Before proceeding to the explicit computation of the potential, let's pause to discuss some intricacies in the choice of the Green's function.
\\
\paragraph{A digression on Green's functions.}
\textcolor{black}{In many amplitude-based derivations (see e.g.~\cite{Cristofoli:2019neg}) one introduces the two-body Green
function
\begin{equation}
  \mathcal{G}(\boldsymbol{k},\boldsymbol{\ell})
  =\frac{1}{E_{\boldsymbol{k}}-E_{\boldsymbol{\ell}}+i0},
  \qquad
  E_{\boldsymbol{\ell}}=\omega_1(\boldsymbol{\ell})+\omega_2(\boldsymbol{\ell})\, .
\end{equation}
In the main text, we instead employ the representation~\eqref{4.4g}, whose leading singular term comes
with the opposite $i0$-prescription. This apparent mismatch is harmless for the conservative observables
we compute: the two choices become equivalent inside the loop integrals once one is allowed to add
terms that are \emph{regular} at $\boldsymbol{\ell}^{2}=\boldsymbol{k}^{2}$ (the ``regular terms'' in~\eqref{4.4g}).
To see this explicitly, Taylor-expand $\mathcal{G}(\boldsymbol{k},\boldsymbol{\ell})$ around
$\boldsymbol{\ell}\simeq\boldsymbol{k}$. Writing the energy difference as an expansion in
$\Delta\equiv(\boldsymbol{\ell}^{2}-\boldsymbol{k}^{2})$, one finds}
\textcolor{black}{\begin{align}
  E_{\boldsymbol{k}}-E_{\boldsymbol{\ell}}+i0
  &=
  -\left.\frac{\partial E_{\boldsymbol{\ell}}}{\partial \boldsymbol{\ell}^{2}}\right|_{\boldsymbol{\ell}^{2}=\boldsymbol{k}^{2}} \Delta
  -\frac12\left.\frac{\partial}{\partial \boldsymbol{\ell}^{2}}
  \Big(\frac{\partial E_{\boldsymbol{\ell}}}{\partial \boldsymbol{\ell}^{2}}\Big)\right|_{\boldsymbol{\ell}^{2}=\boldsymbol{k}^{2}}\Delta^{2}
  +\cdots + i0\,, \nonumber\\[2pt]
  \left.\frac{\partial E_{\boldsymbol{\ell}}}{\partial \boldsymbol{\ell}^{2}}\right|_{\boldsymbol{\ell}^{2}=\boldsymbol{k}^{2}}
  &=\frac{1}{2\,\xi(\boldsymbol{k})\,E_{\boldsymbol{k}}}\,,\qquad
  \left.\frac{\partial}{\partial \boldsymbol{\ell}^{2}}
  \Big(\frac{\partial E_{\boldsymbol{\ell}}}{\partial \boldsymbol{\ell}^{2}}\Big)\right|_{\boldsymbol{\ell}^{2}=\boldsymbol{k}^{2}}
  =-\frac{1-3\xi(\boldsymbol{k})}{4\,\xi(\boldsymbol{k})^{3}\,E_{\boldsymbol{k}}^{3}}\, .
\end{align}}
\textcolor{black}{
In the large center-of-mass energy regime (large $E_{\boldsymbol{k}}$) this yields
\begin{equation}
  \mathcal{G}(\boldsymbol{k},\boldsymbol{\ell})
  =-\frac{2E_{\boldsymbol{k}}\,\xi(\boldsymbol{k})}{\boldsymbol{\ell}^{2}-\boldsymbol{k}^{2}+i0}
  +\frac{3\xi(\boldsymbol{k})-1}{2\,\xi(\boldsymbol{k})\,E_{\boldsymbol{k}}}
  +\mathcal{O}\!\left(\frac{\boldsymbol{\ell}^{2}-\boldsymbol{k}^{2}}{E_{\boldsymbol{k}}^{3}}\right),
  \qquad
  \xi(\boldsymbol{k})=\frac{\omega_1(\boldsymbol{k})\omega_2(\boldsymbol{k})}{(\omega_1(\boldsymbol{k})+\omega_2(\boldsymbol{k}))^{2}}\, .
  \label{eq:Gexp-large-s}
\end{equation}
On the other hand, our choice~\eqref{4.4g} has the leading behavior}
\textcolor{black}{
\begin{equation}
  G(\boldsymbol{k},\boldsymbol{\ell})
  =\frac{1}{(2\pi)^{3}}\frac{1}{2\sqrt{s}}\,
  \frac{1}{\boldsymbol{\ell}^{2}-\boldsymbol{k}^{2}-i0}+\cdots
  \;\;\longrightarrow\;\;
  \frac{1}{(2\pi)^{3}}\frac{1}{2E_{\boldsymbol{k}}}\,
  \frac{1}{\boldsymbol{\ell}^{2}-\boldsymbol{k}^{2}-i0}+\cdots\, ,
  \label{eq:G-ours-leading}
\end{equation}
which differs from~\eqref{eq:Gexp-large-s} (up to normalization) by an overall sign and by the sign of the
$i0$-prescription. Nevertheless, both prescriptions lead to the same integrated result for the class of
amplitudes we need, namely
\begin{equation}
  \mathcal{A}=\int\!\frac{d^{d}\boldsymbol{\ell}}{(2\pi)^{d}}\,
  \frac{\mathrm{GF}(\boldsymbol{k},\boldsymbol{\ell})}{|\boldsymbol{\ell}-\boldsymbol{k}|^{2}\,|\boldsymbol{\ell}-\boldsymbol{k}'|^{2}}\, .
\end{equation}
Indeed, after the standard manipulations, one encounters an integral of the schematic form
\begin{equation}
  \mathcal{A}[G]=
  \int d^{d}\boldsymbol{\ell}\,
  \frac{1}{\big(\boldsymbol{\ell}\!\cdot\!\tilde{\boldsymbol{q}}_{\perp}-i0\big)\,\boldsymbol{\ell}^{2}\,(\boldsymbol{\ell}-\boldsymbol{q})^{2}}\, .
\end{equation}
Now perform the change of variables $\boldsymbol{\ell}\to-\boldsymbol{\ell}$ together with
$\boldsymbol{q}\to-\boldsymbol{q}$ (which leaves the integration domain invariant). This gives
\begin{equation}
  \mathcal{A}[G]=
  -\int d^{d}\boldsymbol{\ell}\,
  \frac{1}{\big(\boldsymbol{\ell}\!\cdot\!\tilde{\boldsymbol{q}}_{\perp}+i0\big)\,\boldsymbol{\ell}^{2}\,(\boldsymbol{\ell}-\boldsymbol{q})^{2}}
  \equiv \mathcal{A}[\mathcal{G}]\, ,
\end{equation}
showing that the two seemingly different Green functions yield the same result once integrated.
Finally, the ellipsis in~\eqref{eq:G-ours-leading} can be fixed by matching to~\eqref{eq:Gexp-large-s} up to
terms analytic at $\boldsymbol{\ell}^{2}=\boldsymbol{k}^{2}$; a convenient choice is
\begin{equation}
  G(\boldsymbol{k},\boldsymbol{\ell})=\frac{1}{(2\pi)^{3}}
  \left[
    \frac{1}{2E_{\boldsymbol{k}}(\boldsymbol{\ell}^{2}-\boldsymbol{k}^{2}-i0)}
    +\frac{3\xi(\boldsymbol{k})-1}{8\,\xi^2(\boldsymbol{k})\,E_{\boldsymbol{k}}^3}
    +\mathcal{O}\!\left(\frac{\boldsymbol{\ell}^{2}-\boldsymbol{k}^{2}}{E_{\boldsymbol{k}}^{5}}\right)
  \right],
\end{equation}
where the subleading pieces are precisely the ``regular terms'' alluded to below~\eqref{4.4g}.}
\\\\
Therefore, the potential in position space can be computed by taking the Fourier transform of \eqref{4.5h}. We have the one-loop and tree-level amplitude. Then our first task to compute the iterative amplitude  \eqref{4.5h}, which yields (we now work in $d$ dimension to isolate the IR divergence),
\begin{align}
  \mathbfcal{A}_{\texttt{Iterative}}=  \int d^d\boldsymbol{\ell}\,\mathbfcal{A}_{\texttt{tree}}(\boldsymbol{k},\boldsymbol{\ell}) G(\boldsymbol{k},\boldsymbol{\ell})\mathbfcal{A}_{\texttt{tree}}(\boldsymbol{\ell},\boldsymbol{k'})
\end{align}
where the tree amplitude is given by,
\begin{align}
\begin{split}
    \mathbfcal{A}_{\texttt{tree}}(\boldsymbol{k},\boldsymbol{\ell})=&-\frac{1}{|\boldsymbol{\ell}-\boldsymbol{k}|^2}\left(-\frac{4  a^2 m_1 m_2}{m_p^2}+16  \alpha _1 \alpha _2 \texttt{g}_ce^2 m_1 m_2 \sigma  +\frac{ m_1^2 m_2^2 \left(-2 \sigma ^2 +1\right)}{2 m_p^2}\right)\,,\\ &
    +\left(4  \alpha _1 \alpha _2 \texttt{g}_c e^2 -\frac{ 2 m_2 m_1 \sigma   }{4  m_p^2}\right)\,,\\ &
    =-c_{1}(\sigma)\frac{1}{|\boldsymbol{\ell}-\boldsymbol{k}|^2}+c_2(\sigma)\,.
    \end{split}
    \end{align}
Therefore, the iterative Born-subtracted amplitude is given by,
\begin{align}
    \begin{split}
        \mathbfcal{A}_{\texttt{Iterative}}=\mathbfcal{A}^{(1)}_{\texttt{Iterative}}+\mathbfcal{A}^{(2)}_{\texttt{Iterative}}+\mathbfcal{A}^{(3)}_{\texttt{Iterative}}
    \end{split}
\end{align}
where,
\begin{align}
    \begin{split}\mathbfcal{A}^{(1)}_{\texttt{Iterative}}(\boldsymbol{k},\boldsymbol{k'})&= \frac{c_1^2(\sigma)}{2E_{\boldsymbol{k}}} \int \frac{d^d\boldsymbol{\ell}}{(2\pi)^3}\,\frac{1}{|\boldsymbol{\ell}-\boldsymbol{k}|^2(\boldsymbol{\ell}^2-\boldsymbol{k}^2-i0)|\boldsymbol{\ell}-\boldsymbol{k'}|^2}\,,\\ &
    =\frac{c_1^2(\sigma)}{2E_{\boldsymbol{k}}}  \int \frac{d^d\boldsymbol{\ell}}{(2\pi)^3}\frac{1}{(\boldsymbol{\ell}^2+2\boldsymbol{\ell}\cdot \boldsymbol{k}-i0)\boldsymbol{\ell}^2(\boldsymbol{\ell}-\boldsymbol{q})^2}\,.\label{4.20jj}
    \end{split}
\end{align}
As before, the integral can be done using the soft loop expansion.
\begin{align}
    \begin{split}\int_{\boldsymbol{\ell\sim\boldsymbol{q}}} \frac{d^d\boldsymbol{\ell}}{(2\pi)^3}\frac{1}{(\boldsymbol{\ell}^2+2\boldsymbol{\ell}\cdot \boldsymbol{k}-i0)\boldsymbol{\ell}^2(\boldsymbol{\ell}-\boldsymbol{q})^2}&=\int_{\boldsymbol{\ell\sim\boldsymbol{q}}} \frac{d^d\boldsymbol{\ell}}{(2\pi)^3}\frac{1}{(\boldsymbol{\ell}^2+2\boldsymbol{\ell}\cdot (\frac{\tilde{\boldsymbol{q}}_{\perp}-\boldsymbol{q}}{2})-i0)\boldsymbol{\ell}^2(\boldsymbol{\ell}-\boldsymbol{q})^2}\,,\\ &    \hspace{0cm}=\int_{\boldsymbol{\ell\sim\boldsymbol{q}}} \frac{d^d\boldsymbol{\ell}}{(2\pi)^3}\left(1-\frac{\boldsymbol{\ell}^2-\boldsymbol{\ell}\cdot \boldsymbol{q}}{\boldsymbol{\ell}\cdot\tilde{\boldsymbol{q}}_{\perp}}+\cdots\right) \frac{1}{(\boldsymbol{\ell}\cdot\boldsymbol{q_{\perp}}-i0)\boldsymbol{\ell}^2(\boldsymbol{\ell-q})^2}\label{4.10o}
    \end{split}
\end{align}
where, $\tilde{\boldsymbol{q}}_{\perp}=\boldsymbol{q}+2\boldsymbol{k}$ and also note that $\tilde{\boldsymbol{q}}_{\perp}\cdot \boldsymbol{q}=0$. Keeping this in mind the first term in \eqref{4.10o} takes the form,
\begin{align}
    \textstyle{\mathbfcal{A}^{(1)}_{\texttt{Iterative}}(\boldsymbol{k},\boldsymbol{k'})=\frac{c^2_1(\sigma)}{2E_{\boldsymbol{k}}}\int_{\boldsymbol{\ell\sim\boldsymbol{q}}} \frac{d^d\boldsymbol{\ell}}{(2\pi)^3} \frac{1}{(\boldsymbol{\ell}\cdot\tilde{\boldsymbol{q}}_{\perp}-i0)\boldsymbol{\ell}^2(\boldsymbol{\ell-q})^2}}&\xrightarrow[\boldsymbol{q}\to -\boldsymbol{q}]{\boldsymbol{\ell}\to -\boldsymbol{\ell}}-\frac{c^2_{1}(\sigma)}{2E_k}   \int d^{d}\boldsymbol{\ell}\frac{1}{(\boldsymbol{\ell}\cdot \tilde{\boldsymbol{q}}_{\perp}+i0)\boldsymbol{\ell}^2(\boldsymbol{\ell}-\boldsymbol{q})^2},
\end{align}
and adding the two, we get,
\begin{align}
   \begin{split}\mathbfcal{A}^{(1)}_{\texttt{Iterative}}(\boldsymbol{k},\boldsymbol{k'})= \frac{\pi i c_1^2(\sigma)}{2E_{\boldsymbol{k}}}\int \frac{d^d\boldsymbol{\ell}}{(2\pi)^3} \frac{\delta(\boldsymbol{\ell}\cdot\tilde{\boldsymbol{q}}_{\perp})}{\boldsymbol{\ell}^2(\boldsymbol{\ell}-\boldsymbol{q})^2}&=\frac{c_1^2(\sigma)}{2E_{\boldsymbol{k}}}\frac{i\pi}{2\pi |\tilde{\boldsymbol{q}}_{\perp}|}\int \frac{d^{d-1}}{(2\pi)^2}\boldsymbol{\ell} \frac{1}{\boldsymbol{\ell}^2(\boldsymbol{\ell}-\boldsymbol{q})^2}\,,\\ &
 =\frac{c_1^2(\sigma)}{2E_{\boldsymbol{k}}}\frac{i}{2|\boldsymbol{q}_\perp|}\left(\frac{ \log (|\boldsymbol{q}|^2)+\gamma_{E} -\log (4 \pi )}{2 \pi  |\boldsymbol{q}|^2}-\frac{1}{2 \pi |\boldsymbol{ q}|^2 \epsilon }\right)\,.
    \end{split}
\end{align}
Again note that, in the soft limit $|\boldsymbol{q}_\perp|=\sqrt{4\boldsymbol{k}^2-\boldsymbol{q}^2}\sim {2}|\boldsymbol{k}|$.
As we immediately see that the iterative amplitude has an IR divergent piece which is given by,
\begin{align}
    \mathbfcal{A}_{\texttt{Iterative}}^{(1)}(\boldsymbol{k},\boldsymbol{k'})\Bigg|_{\texttt{IR div.}}= -\frac{i c_1^2(\sigma)}{16 \pi E_{\boldsymbol{k}}|\boldsymbol{k}||\boldsymbol{q}|^2\,\epsilon}=-\frac{ic_1^2(\sigma)}{16\pi m_1 m_2 \sqrt{\sigma^2-1}|\boldsymbol{q}|^2\epsilon}\,.
\end{align}
On the other hand,  the one-loop amplitude computed in \eqref{3.54j} has an IR divergent piece,
\begin{align}
    \begin{split}\mathbfcal{A}_{\texttt{1-loop}}\Bigg|_{\texttt{IR div.}}&=-\frac{i m_1 m_2 \left(8 \left(a^2-4 \alpha _1 \alpha _2 e^2 \sigma  \texttt{g}_c m_p^2\right)+m_1 m_2 \left(2 \sigma ^2-1\right)\right){}^2}{64 \pi  \sqrt{\sigma ^2-1} \epsilon \, m_p^4 |\boldsymbol{q}|^2}\,,\\ & =\frac{-ic_1^2(\sigma)}{16\pi m_1 m_2\sqrt{\sigma^2-1}|\boldsymbol{q}|^2\epsilon}\to \mathbfcal{A}_{\texttt{Iterative}}^{(1)}(\boldsymbol{k},\boldsymbol{k'})\Bigg|_{\texttt{IR div.}} \,.
    \end{split}
\end{align}
From the above analysis, \textit{We observe a complete cancellation of infrared singularities: the IR divergence present in the one-loop amplitude is exactly removed by the iterative (Born) contribution. Hence, the Born (IR)-subtracted Post-Minkowskian potential is IR finite. \textcolor{black}{The same subtraction scheme has been first applied for general relativity in~\cite{Cristofoli:2019neg}, where the Born-subtracted potential was also demonstrated to be IR finite.}}   As we will see, the other integrals appearing in the iterative amplitude are free of IR divergence.
Similarly, we have,
\begin{align}
    \begin{split}
        \mathbfcal{A}_{\texttt{Iterative}}^{(2)}(\boldsymbol{k},\boldsymbol{k'})&\sim-c_1(\sigma)c_2(\sigma)\int d^d\boldsymbol{\ell}\frac{1}{|\boldsymbol{\ell}-\boldsymbol{k}|^2(\boldsymbol{\ell}^2-\boldsymbol{k}^2)}\,,\\ &
        =-c_1(\sigma)c_2(\sigma)\int d^d\boldsymbol{\ell}\frac{1}{\boldsymbol{\ell}^2(\boldsymbol{\ell}^2+2\boldsymbol{\ell}\cdot\boldsymbol{k}-i0)}=0\,\, (\texttt{scaleless in soft limit})\,,\\ &
        \to \mathbfcal{A}_{\texttt{Iterative}}^{(3)}(\boldsymbol{k},\boldsymbol{k'}) \,\,(\texttt{in soft limit})\,.
    \end{split}
\end{align}
However, we are not completely done yet. We see that the iterative Born subtraction only has a superclassical contribution. However, we still have the freedom to add a regular term coming from the Taylor expansion mentioned in the footnote. The Green function takes the form,
\begin{align}
    G(\boldsymbol{k},\boldsymbol{\ell})=\frac{1}{(2\pi)^3}\frac{1}{2\sqrt{s}}\frac{1}{|\boldsymbol{\ell}|^2-|\boldsymbol{k}|^2-i0}+\frac{3\xi-1}{64 \pi^3 \xi^2 {s^{\frac{3}{2}}}}+\cdots\,.
\end{align}
Therefore, the only surviving correction to the iterative amplitude in the soft limit is the following,
\begin{align}
    \begin{split}
        \mathbfcal{A}_{\texttt{Iterative}}'(\boldsymbol{k},\boldsymbol{k'})=c_1^2(\sigma)\frac{3\xi-1}{64\pi^3\xi^2s^{\frac{3}{2}}}\int \frac{d^d\boldsymbol{\ell}}{\boldsymbol{\ell}^2(\boldsymbol{\ell}-\boldsymbol{q})^2}=c_1^2(\sigma)\frac{3\xi-1}{64\xi^2s^{\frac{3}{2}}|\boldsymbol{q}|}\,.
    \end{split}
\end{align}
Therefore, the one-loop classical potential in momentum space is given by,
\begin{align}
V(\boldsymbol{k},|\boldsymbol{q}|)=\mathbfcal{A}_{\texttt{1-loop}}(\boldsymbol{k},|\boldsymbol{q}|)\Big|_{\frac{1}{|\boldsymbol{q}|}}-\frac{c_1^2(\sigma)(3\xi-1)}{64\xi^2 E_{\boldsymbol{k}}^3 |\boldsymbol{q}| }\,.
\end{align}
The two-body classical  potential can be computed by taking the Fourier transform of the momentum space amplitude,
\begin{align}   
\begin{split}
V_{\texttt{1-loop}}(\boldsymbol{k},|\boldsymbol{r}|)=-\frac{1}{
4m_1m_2}\int \frac{d^3\boldsymbol{q}}{(2\pi)^3}e^{i\boldsymbol{q}\cdot \boldsymbol{r}}\,V(\boldsymbol{k},|\boldsymbol{q}|) =-\frac{\mathcal{C}_2(\sigma)}{8m_1m_2\pi^2 r^2}\,.
\end{split}
\end{align}
where,
% \begin{align}
%     \begin{split}
%    \mathcal{C}_{2}(\sigma)&=     \frac{2 \sigma  \left(\sqrt{\sigma ^2-1}-\sigma \right)+1}{512 \left(\sigma -\sqrt{\sigma ^2-1}\right)^2 m_p^4}\Bigg[-128 m_2 \left(a^4+16 \alpha _1^2 \alpha _2^2 e^4 \texttt{g}_c^2 m_p^4\right)+m_2 m_1^2 \Big(4 a^2 \left(\sigma ^2-17\right)\\&+32 \alpha _2 e^2 \left(\alpha _2 \left(3 \sigma ^2-1\right)+12 \alpha _1 \sigma \right) \texttt{g}_c m_p^2+m_2^2 \left(3-15 \sigma ^2\right)\Big)+4 m_1 \Bigg(-32 \left(a^4+16 \alpha _1^2 \alpha _2^2 e^4 \texttt{g}_c^2 m_p^4\right)\\ &+m_2^2 \left(a^2 \left(\sigma ^2-17\right)+8 \alpha _1 e^2 \left(\alpha _1 \left(3 \sigma ^2-1\right)+12 \alpha _2 \sigma \right) \texttt{g}_c m_p^2\right)\\ &+32 a m_2 \left(4 a b+\sqrt{2} e^2 \left(-4 \alpha _2 \alpha _1 \sigma +\alpha _1^2+\alpha _2^2\right) \texttt{g}_c m_p^2\right)\Bigg)+3 m_2^2 m_1^3 \left(1-5 \sigma ^2\right)\Bigg]-\frac{c_1^2(3\xi-1)}{32\pi^3\xi^2 E_{\boldsymbol{k}}^3} ,\,.\label{4.21h}
%     \end{split}
% \end{align}
\begin{align}
    \begin{split}
         \mathcal{C}_{2}(\sigma) &= \frac{a^4 \left(m_1+m_2\right)}{4 m_p^4}-\frac{a^2 m_1 m_2 \left(128 b+\left(m_1+m_2\right) \left(\sigma ^2-17\right)\right)}{128 m_p^4}\\
         & -\frac{\sqrt{2} a e^2 m_1 m_2 \left(-4 \alpha _2 \alpha _1 \sigma +\alpha _1^2+\alpha _2^2\right) \texttt{g}_c}{4 m_p^2}  +\frac{3 m_1^2 m_2^2 \left(m_1+m_2\right) \left(5 \sigma ^2-1\right)}{512 m_p^4} \\
         &+ 4 \alpha _1^2 \alpha _2^2 e^4 \left(m_1+m_2\right) \texttt{g}_c^2 -\frac{e^2 m_1 m_2 \texttt{g}_c \left(\left(3 \sigma ^2-1\right) \left(\alpha _1^2 m_2+\alpha _2^2 m_1\right)+12 \alpha _1 \alpha _2 \left(m_1+m_2\right) \sigma \right)}{16 m_p^2}\\
         & -\frac{c_1^2(\sigma)(3\xi-1)}{64\xi^2 E_{\boldsymbol{k}}^3  } \,,\label{4.21g}
    \end{split}
\end{align}
with{\footnote{where $a,b$ has mass dimension 1.}} $\sigma=\frac{k_1\cdot k_2}{m_1m_2}\,$. If one turns off the dilaton field and takes the static limit, then the potential written in \eqref{4.21g} agrees with \cite{Bjerrum-Bohr:2002aqa,Faller:2007sy}. Before closing this section, we briefly comment on a mild ambiguity that arises from the freedom to add \emph{regular} (analytic) terms to the Green's function, as mentioned earlier.\\\\
\textcolor{black}{
\textbf{Comment on adding regular pieces in the Green's function.}  The freedom to add \emph{regular} (analytic) terms to the Green’s function only affects the off-shell extension of the integrand and can at most generate
shorter-range contributions to the conservative potential. In particular, such terms do \emph{not} modify the IR-singular piece that is removed by the Born/iteration subtraction.
To see this explicitly, note that adding analytic regular terms produces (schematically) loop integrals of the form
\begin{align}
\mathcal{I}_n(\boldsymbol{q},\boldsymbol{q}_\perp)
\;\equiv\;
\int d^d\boldsymbol{\ell}\;
\frac{\bigl(2\,\boldsymbol{\ell}\!\cdot\!\boldsymbol{q}_\perp\bigr)^n}{\boldsymbol{\ell}^{\,2}\,(\boldsymbol{\ell}-\boldsymbol{q})^{2}}\,,
\qquad n\ge 1,
\label{eq:In_def_refreply}
\end{align}
where $\boldsymbol{q}_\perp$ is transverse, $\boldsymbol{q}\!\cdot\!\boldsymbol{q}_\perp=0$ (mass on-shell condition: $|\boldsymbol{k}|^2=|\boldsymbol{k}+\boldsymbol{q}|^2$). By symmetry, $\mathcal{I}_n$ vanishes identically for odd $n$.
For even $n$, a standard IBP reduction implies that the numerator can only generate powers of
$\boldsymbol{q}^2$ and $\boldsymbol{q}_\perp^{\,2}$, so that one may write
\begin{align}
\mathcal{I}_n(\boldsymbol{q},\boldsymbol{q}_\perp)
=
C_n(d)\,(\boldsymbol{q}^{2})^{n/2}(\boldsymbol{q}_\perp^{\,2})^{n/2}
\int d^d\boldsymbol{\ell}\,\frac{1}{\boldsymbol{\ell}^{\,2}(\boldsymbol{\ell}-\boldsymbol{q})^{2}}
\;\propto\;
(\boldsymbol{q}^2)^{\frac{n-1}{2}}
\qquad (\text{even } n),
\label{eq:In_scaling_refreply}
\end{align}
where in the last step we used the well-known scaling of the massless bubble integral in $d=3-2\epsilon$ spatial dimensions,
$\int d^d\boldsymbol{\ell}\, [\boldsymbol{\ell}^2(\boldsymbol{\ell}-\boldsymbol{q})^2]^{-1}\propto (\boldsymbol{q}^2)^{d/2-2}\sim (\boldsymbol{q}^2)^{-1/2-\epsilon}$.
Crucially, for $n\ge 1$ the behavior \eqref{eq:In_scaling_refreply} is \emph{less singular} in the soft limit than the $n=0$ term, and therefore it does not contribute to the IR divergence that is cancelled by the iteration (Born) subtraction. Finally, the corresponding contribution to the conservative potential is obtained by Fourier transform.
Using $\int d^3\boldsymbol{q}\,e^{i\boldsymbol{q}\cdot\boldsymbol{r}}\,|\boldsymbol{q}|^{\,m}\sim |r|^{-m-3}$ (up to constants), one finds
\begin{align}
V_n(r)\;\sim\;\int d^3\boldsymbol{q}\;e^{i\boldsymbol{q}\cdot\boldsymbol{r}}\,
(\boldsymbol{q}^2)^{\frac{n-1}{2}}
\;\propto\;\frac{1}{|\boldsymbol{r}|^{\,n+2}}\,, \qquad n=2,4,6,\cdots,
\label{eq:Vn_falloff_refreply}
\end{align}
i.e.\ a contribution that is parametrically shorter-ranged (and more suppressed at large $r$) than the leading long-range terms.
Therefore, purely analytic regular terms cannot alter the IR-subtracted \emph{long-range} potential: their effect is confined to short-distance/contact contributions and, more generally, to off-shell terms that can be absorbed by local field redefinitions (or equivalently by canonical transformations).
It is useful to recall how the relevant non-analytic structures in the momentum-transfer expansion map into the large-distance behavior of the potential at $\mathcal{O}(G^2)$.
In three spatial dimensions, the Fourier transform implies the schematic correspondences
\begin{align}
\frac{1}{|\boldsymbol{q}|^{2}}\;\longleftrightarrow\;\frac{G^2}{r\hbar}\,,\qquad
\frac{1}{|\boldsymbol{q}|}\;\longleftrightarrow\;\frac{G^2 \hbar^0}{r^{2}}\,,\qquad
\log(\boldsymbol{q}^{2})\;\longleftrightarrow\;\frac{G^2\hbar}{r^{3}}\,,
\end{align}
which are naturally interpreted as the super-classical, classical, and quantum contributions, respectively.
In our analysis, the apparent super-classical ($\sim 1/\hbar$) pole is spurious and is removed by the standard leading Born subtraction, leaving an IR-finite long-range result.
Regular terms, on the other hand, only introduce an ambiguity in the definition of the post-Minkowskian potential at the level of off-shell completion.
As emphasized in \cite{Correia:2024jgr,Correia:2024yfx}, different choices of regular terms lead to potentials that differ by off-shell pieces, and hence they do not affect on-shell, gauge-invariant observables (such as the scattering angle or other conservative observables derived from the S-matrix).  }
Next, we discuss eikonal exponentiation in momentum space and deduce the scattering angle.

\section{Eikonal exponentiation and the scattering angle}\label{ch4:sec5}
In this section, we will study the eikonal exponentiation of the conservative scattering amplitude directly in the momentum space following \cite{Parra-Martinez:2020dzs}. Traditionally, in eikonal exponentiation, it is well known that the eikonal phase can be computed by taking the Fourier transform of the scattering amplitude in impact-parameter space.
The primary reason to do the computation in momentum space is to avoid additional IR divergence coming from the tree-level amplitudes arising from the Coulomb interaction. However, one must pay the computational cost in momentum space: simple products in impact parameter space give rise to convolutions in momentum space, which can be computed using the iterative bubble integrals. The statement of eikonal exponentiation reduces to the following: one can write the scattering amplitude as a \textit{convolutional exponential} of the eikonal phase,
\begin{align}
\begin{split}
    i\mathbfcal{A}(\sigma,-q^2)&=\textrm{cexp}\left(i\delta(\sigma,-q^2)\right)-1\,,\\ &
    =i\delta(\sigma,-q^2)-\frac{1}{2!}\delta(\sigma,-q^2)\otimes \delta(\sigma,-q^2)-i\frac{1}{3!}\delta(\sigma,-q^2)\otimes \delta(\sigma,-q^2)\otimes\delta(\sigma,-q^2)+\cdots
    \end{split}
\end{align}
where the convolution is defined as an integral over the $d-2$ dimensional transverse space.
\begin{align}    \varpi_1(\boldsymbol{q}_{\perp})\otimes \varpi_2(\boldsymbol{q}_{\perp}):=\frac{1}{N}\int d\mu_{\boldsymbol{\ell}_{\perp}}^{d-2}\,\varpi_1(\boldsymbol{q}_{\perp})\varpi_2(\boldsymbol{q}_{\perp}-\boldsymbol{\ell}_{\perp})\,,
\end{align}
where $N=4 m_1 m_2 \sqrt{\sigma^2-1}$ is the normalization factor. Similarly, one can write the inverse relation as,
\begin{align}
    \begin{split}
        \delta(\sigma,-q^2)&=-i\textrm{clog}(1+i \mathbfcal{A})\,, \\
        &=\mathbfcal{A}(\sigma ,-q^2)-\frac{i}{2}\mathbfcal{A}(\sigma ,-q^2) \otimes\mathbfcal{A}(\sigma ,-q^2) -\frac{1}{3}\mathbfcal{A}(\sigma ,-q^2) \otimes \mathbfcal{A}(\sigma ,-q^2) \otimes\mathbfcal{A}(\sigma ,-q^2)+\cdots\,.
    \end{split}
\end{align}
This can be written as 
\begin{align}
    \begin{split}
        \delta= \delta^{(0)}+\delta^{(1)}+\delta^{(2)}+\cdots
    \end{split}
\end{align}
where $ \delta^{(L)}$ has $\mathcal{O}(e^{2L+2})$,  $\mathcal{O}(G^{L+1})$ and $\mathcal{O}(e^{2L} G^{L})$ terms and so on hence we can write phases as
\begin{align}
    \begin{split}
    &\delta^{(0)}=\mathbfcal{A}^{\textrm{tree}}\,, \\
    & \delta^{(1)}=\mathbfcal{A}^{\textrm{1-loop}} -\frac{i}{2} \mathbfcal{A}^{\textrm{tree}} \otimes \mathbfcal{A}^{\textrm{tree}}\,.
    \end{split}
\end{align}
So we have
\begin{align}
    \delta^{(0)}(\sigma, -q^2)= \frac{1}{q^2} \left( -\frac{4 a^2 m_1 m_2}{m_p^2} + 16 \alpha _1 \alpha _2 e^2   m_1 m_2 \sigma   \texttt{g}_c+\frac{m_1^2 m_2^2 \left(-2 \sigma ^2+2 \sigma ^2 \epsilon +1\right)}{(2-2 \epsilon ) m_p^2} \right) \label{del0p}
\end{align}
where the $\mathcal{O}(-q^0)$ term of $\mathbfcal{A^{\textrm{tree}}}$ won't contribute as they lead to $\delta(\boldsymbol{b_e})$ (which is simply zero as $\boldsymbol{b_e} \neq 0$) when fourier transformed to impact parameter space.
Now using 
\begin{align}
    \frac{1}{-q^2}\otimes\frac{1}{-q^2}=\frac{1}{\boldsymbol{q^2_\perp}}\otimes\frac{1}{\boldsymbol{q^2_\perp}}=\frac{1}{N} (4 \pi )^{\epsilon -1}  \frac{  \Gamma (-\epsilon )^2 \Gamma (\epsilon +1)}{\Gamma (-2 \epsilon )} (\boldsymbol{q}^2_\perp)^{-\epsilon -1}
\end{align}
we get the following, 
\begin{align}
    \begin{split}
        &\delta^{(1)}(\sigma, -q^2) = \frac{i \left(2 m_2 m_1 \sigma +m_1^2+m_2^2\right) \left(64 a^4+m_1^2 m_2^2 (1-2 \sigma ^2\right)^2)}{128 \pi  m_1 m_2 \left(\sigma ^2-1\right)^{3/2} m_p^4} \log{(\boldsymbol{q}^2_\perp )}  \\
        & + \frac{\left(m_1+m_2\right) \left(128 a^4-4 a^2 m_1 m_2 \left(\sigma ^2-17\right)+3 m_1^2 m_2^2 \left(5 \sigma ^2-1\right)\right)-512 a^2 b m_1 m_2}{512 m_p^4 |\boldsymbol{q}_\perp|} \\
        & -\frac{e^2 m_1 m_2 \texttt{g}_c \left(4 \sqrt{2} a \left(\alpha _1^2+\alpha _2^2\right)+4 \alpha _1 \alpha _2 \sigma  \left(3 \left(m_1+m_2\right)-4 \sqrt{2} a\right)+\left(3 \sigma ^2-1\right) \left(\alpha _1^2 m_2+\alpha _2^2 m_1\right)\right)}{16 m_p^2 |\boldsymbol{q}_\perp|} \\
        &  +\frac{4 \alpha _1^2 \alpha _2^2 e^4 \left(m_1+m_2\right) \texttt{g}_c^2}{|\boldsymbol{q_\perp}|}\,.
        \label{del1p}
    \end{split}
\end{align}
Note that the $\mathcal{O}(|\boldsymbol{q}_\perp|^{-2}) $ term of the 1-loop amplitude is completely canceled at $\epsilon^0$ order by the tree convolution. We also note that the $\delta^{(1)}$ has an imaginary part. However, the imaginary part of the eikonal phase does not contribute to the conservative scattering angle. Physically, the imaginary part of the eikonal phase measures the number of massless particles emitted from the two-body system. We now proceed to the computation of the eikonal phase in impact-parameter space by taking the Fourier transform. 
\begin{align}
    \delta(\sigma, \boldsymbol{b}_e )=\frac{1}{N} \int\frac{d^{D-2}\boldsymbol{q}_\perp}{(2 \pi)^{D-2}} e^{i \boldsymbol{b}_e \cdot \boldsymbol{q}_\perp} \delta(\sigma,\boldsymbol{q}_\perp)\,.
\end{align}
So using \eqref{del0p} and \eqref{del1p} respectively we have
\begin{flalign}
   \hspace{0cm}  \delta^{(0)}(\sigma, \boldsymbol{b}_e )&= \left( -\frac{a^2}{4 \pi  \sqrt{\sigma ^2-1} m_p^2}+\frac{\alpha _1 \alpha _2 e^2 \sigma  \texttt{g}_c}{\pi  \sqrt{\sigma ^2-1}}-\frac{m_1 m_2 \left(2 \sigma ^2-1\right)}{32 \pi  \sqrt{\sigma ^2-1} m_p^2} \right) \log(\boldsymbol{b}^2_e)&
\end{flalign}
and,
\begin{align}
    \begin{split}
        \delta^{(1)}(\sigma, \boldsymbol{b}_e ) &= \frac{128 a^4 \left(m_1+m_2\right)-4 a^2 m_1 m_2 \left(128 b+\left(m_1+m_2\right) \left(\sigma ^2-17\right)\right)+3 m_1^2 m_2^2 \left(m_1+m_2\right) \left(5 \sigma ^2-1\right)}{4096 \pi  m_1 m_2 \sqrt{\sigma ^2-1}  m_p^4 |\boldsymbol{b}_e| } \\ & -\frac{e^2 \texttt{g}_c \left(4 \sqrt{2} a \left(\alpha _1^2+\alpha _2^2\right)+4 \alpha _1 \alpha _2 \sigma  \left(-4 \sqrt{2} a+3 m_1+3 m_2\right)+\left(3 \sigma ^2-1\right) \left(\alpha _1^2 m_2+\alpha _2^2 m_1\right)\right)}{128 \pi  \sqrt{\sigma ^2-1} m_p^2 |\boldsymbol{b}_e|} \\
        & + \frac{\alpha _1^2 \alpha _2^2 e^4 \left(m_1+m_2\right) \texttt{g}_c^2}{2 \pi  m_1 m_2 \sqrt{\sigma ^2-1} |\boldsymbol{b}_e|}\,.
    \end{split}
\end{align}
\textbf{\textit{Scattering angle:}}\\
The phase in impact parameter space and the amplitude in momentum space are related as
\begin{align}
    \mathbfcal{A}(\sigma,-q^2)=\int d^{D-2} \boldsymbol{b}_e \left( e^{i \delta(\sigma,\boldsymbol{b}_e)} -1 \right) e^{-i \boldsymbol{b}_e \cdot \boldsymbol{q}} 
\end{align}
and stationary phase approximation yields the relation
\begin{align}
    \boldsymbol{q}=-\frac{\partial}{\partial\boldsymbol{b}_e} \delta(\sigma,\boldsymbol{b}_e)\,. \label{5.14}
\end{align}
In the centre-of-mass frame, $\boldsymbol{q}$, scattering angle $\chi$ and the COM momentum $\boldsymbol{p}$ are related as 
\begin{align}
    | \boldsymbol{q}|=2 |\boldsymbol{k}| \sin{\frac{\chi}{2}} \label{5.15}
\end{align}
where in terms of centre-of-mass energy
\begin{align}
    |\boldsymbol{k}|=\frac{m_1 m_2 \sqrt{\sigma^2 - 1}}{E} = \frac{m_1 m_2 \sqrt{\sigma^2 - 1}}{\sqrt{m_1^2 + m_2^2 +2 m_1 m_2 \sigma}}\,.
\end{align}
So from \eqref{5.14} and \eqref{5.15} we have
\begin{align}
    \sin{\frac{\chi}{2}}=-\frac{1}{2 |\boldsymbol{k}|} \frac{\partial}{\partial |\boldsymbol{b}_e|} \Re \delta(\sigma,\boldsymbol{b}_e)\,. \label{5.17}
\end{align}
Now we can write RHS of \eqref{5.17} by order
\begin{flalign}
    \frac{\chi^{(0)}}{2}&=\frac{a^2}{4 \pi  |\boldsymbol{k}||\boldsymbol{b}_e| \sqrt{\sigma ^2-1} m_p^2} - \frac{\alpha _1 \alpha _2 e^2 \sigma  \texttt{g}_c}{\pi  |\boldsymbol{k}||\boldsymbol{b}_e| \sqrt{\sigma ^2-1}}+\frac{m_1 m_2 \left(2 \sigma ^2-1\right)}{32 \pi  |\boldsymbol{k}||\boldsymbol{b}_e| \sqrt{\sigma ^2-1} m_p^2}\,, &
\end{flalign}

\begin{align}
    \begin{split}
        \frac{\chi^{(1)}}{2} &= \frac{a^4 \left(m_1+m_2\right)}{64 \pi  m_1 m_2 |\boldsymbol{k}||\boldsymbol{b}^2_e| \sqrt{\sigma ^2-1} m_p^4}-\frac{a^2 \left(128 b+\left(m_1+m_2\right) \left(\sigma ^2-17\right)\right)}{2048 \pi  |\boldsymbol{k}||\boldsymbol{b}^2_e| \sqrt{\sigma ^2-1} m_p^4}  +\frac{3 m_1 m_2 \left(m_1+m_2\right) \left(5 \sigma ^2-1\right)}{8192 \pi  |\boldsymbol{k}||\boldsymbol{b}^2_e| \sqrt{\sigma ^2-1} m_p^4} \\
        & -\frac{e^2 \texttt{g}_c \left(4 \sqrt{2} a \left(\alpha _1^2+\alpha _2^2\right)+4 \alpha _1 \alpha _2 \sigma  \left(-4 \sqrt{2} a+3 m_1+3 m_2\right)+\left(3 \sigma ^2-1\right) \left(\alpha _1^2 m_2+\alpha _2^2 m_1\right)\right)}{256 \pi  |\boldsymbol{k}||\boldsymbol{b}^2_e| \sqrt{\sigma ^2-1} m_p^2} \\
        & + \frac{\alpha _1^2 \alpha _2^2 e^4 \left(m_1+m_2\right) \texttt{g}_c^2}{4 \pi  m_1 m_2 |\boldsymbol{k}||\boldsymbol{b}^2_e| \sqrt{\sigma ^2-1}}\,.
    \end{split}
\end{align}
The above results can be written in terms of angular momentum 
\begin{align}
    J=|\boldsymbol{b} \times \boldsymbol{k}|,
\end{align}
where $\boldsymbol{b}$ is the impact parameter perpendicular to the incoming centre-of-mass momentum $\boldsymbol{p}$. Note, however, this impact parameter is different from the $\boldsymbol{b}_e$ occurring in the eikonal phase, which points in the direction of momentum transfer $\boldsymbol{q}$. The magnitude of $\boldsymbol{b}_e$ and $\boldsymbol{b}$ are related by
\begin{align}
    |\boldsymbol{b}|=|\boldsymbol{b}_e| \cos{\frac{\chi}{2}}.
\end{align}
This difference is unimportant for small-angle scattering, and it will only matter at order $J^{-3}$. So in terms of $J$ and $G$ ($ m_p=(32 \pi G)^{-1/2}$)  we have,
\begin{flalign}
    \frac{\chi^{(0)}}{2}&=\frac{8 a^2 G}{J \sqrt{\sigma ^2-1}}-\frac{\alpha _1 \alpha _2 e^2 \sigma  \texttt{g}_c}{\pi  J \sqrt{\sigma ^2-1}}+\frac{G m_1 m_2 \left(2 \sigma ^2-1\right)}{J \sqrt{\sigma ^2-1}}\,, &
\end{flalign}

\begin{align}
    \begin{split}
        \frac{\chi^{(1)}}{2}&= \frac{16 \pi  a^4 G^2 \left(m_1+m_2\right)}{J^2 \sqrt{2 m_2 m_1 \sigma +m_1^2+m_2^2}} -\frac{\pi  a^2 G^2 m_1 m_2 \left(128 b+\left(m_1+m_2\right) \left(\sigma ^2-17\right)\right)}{2 J^2 \sqrt{2 m_2 m_1 \sigma +m_1^2+m_2^2}} \\
        & + \frac{3 \pi  G^2 m_1^2 m_2^2 \left(m_1+m_2\right) \left(5 \sigma ^2-1\right)}{8 J^2 \sqrt{2 m_2 m_1 \sigma +m_1^2+m_2^2}} +\frac{\alpha _1^2 \alpha _2^2 e^4 \left(m_1+m_2\right) \texttt{g}_c^2}{4 \pi  J^2 \sqrt{2 m_2 m_1 \sigma +m_1^2+m_2^2}} \\
        & -\frac{e^2 G m_1 m_2 \texttt{g}_c \left(4 \sqrt{2} a \left(\alpha _1^2+\alpha _2^2\right)+4 \alpha _1 \alpha _2 \sigma  \left(-4 \sqrt{2} a+3 m_1+3 m_2\right)+\left(3 \sigma ^2-1\right) \left(\alpha _1^2 m_2+\alpha _2^2 m_1\right)\right)}{8 J^2 \sqrt{2 m_2 m_1 \sigma +m_1^2+m_2^2}}\,.
    \end{split}
\end{align}
If we `turn off' the coupling, i.e $G \rightarrow 0$ \cite{Bern:2021xze} or $e \rightarrow0$ \cite{Mogull:2020sak}, and set the screening constants $a_i \, ,  b_i =0$, we retrieve the known gravitational and (scalar) electromagnetic scattering angles. 
\section{Conclusions and Discussion}
\label{sec:conclusions}
We have analyzed classical scattering of charged, non-spinning compact objects in EMD theory using modern scattering–amplitude methods, with cross–checks based on potential theory. Our goal was to extract the conservative two–body potential, establish the infrared (IR) structure and its cancellation in momentum space, and compute the eikonal phase and associated scattering angle through one loop (classical $2{\rm PM}$ order).


\begin{itemize}
  \item  We have motivated EMD as a low–energy effective description inspired by heterotic string theory, in which additional long–range fields (dilaton and gauge bosons) arise from first principles. This provides a UV-motivated setting in which higher–curvature corrections and extra degrees of freedom naturally coexist, while our concrete computations focused on the leading two-derivative EMD sector.
  \item After fixing conventions and gauges, we derived the propagators and Feynman rules needed for $2\!\to\!2$ scattering. This delivers a set of building blocks separating purely gravitational, electromagnetic, and dilatonic exchanges and their interference.
  \item  Using dimensional regularization, expansion by regions, and IBP reduction, we obtained a basis of master integrals and their soft (classical) expansions that capture the long–range, nonanalytic in momentum transfer relevant for conservative dynamics. These ingredients were assembled into the one–loop amplitudes across the relevant topologies (single–exchange, triangle, box, and penguin).
  \item We computed the momentum–space two–body potential via the Lippmann-Schwinger equation. \textit{A careful EFT/Born subtraction removes iterated long–range exchanges, and we explicitly showed that the resulting momentum–space potential is IR finite, mirroring the situation in GR}. This establishes that the soft amplitude, once dressed by the appropriate iterative counter term, is free of IR divergences at the level relevant for conservative observables.
  \item  The conservative scattering angle was obtained by exponentiating the momentum–space amplitude (eikonal exponentiation) and differentiating the eikonal phase with respect to impact parameter, and the result is given as an explicit closed-form function of gravitational, electromagnetic, and dilatonic couplings. We also verify that, in the appropriate limits, our results for the classical potential and the scattering angle agree precisely with those available in the literature.

\end{itemize}
%\paragraph{Phenomenological implications.}
%\textcolor{black}{Although the present study focuses on hyperbolic encounters (the native regime of PM and eikonal methods), the resulting angle and its impact–parameter dependence feed directly into effective–one–body (EOB) maps and calibrated waveform models. Consequently, our results furnish amplitude–based benchmarks for beyond–GR dynamics and can be propagated into waveform phasing for (i) dynamical capture scenarios, (ii) scattering events in dense environments, and (iii) constraints on dilaton couplings and charge–to–mass ratios using current and next–generation detectors. The IR–finite potential we constructed is particularly well–suited for such mappings.
%\newpage 
\textbf{Future outlook}:
There are several interesting directions to be done. 
\begin{itemize}
  \item \emph{Higher orders.} \textcolor{black}{Extending to two loops ($3{\rm PM}$) would sharpen the conservative sector and test the robustness of IR cancellations beyond one loop in EMD. However, at two-loop and onward radiative corrections appear to emerge. However, we think that, as long as we consider the conservative scattering angle at two loops, the Born subtraction scheme should work. However if we add the radiative corrections (where one needs to compute the boundary integrals both from potential and radiation region) the subtraction scheme may fail. However, in those cases, the EFT prescription has been checked to work perfectly in the context of GR in \cite{Cheung:2018wkq}. We hope to report this soon \cite{arpan2}.}
  \item \emph{Spin and finite size.} Incorporating spins, spin–orbit/spin–spin couplings, and tidal responses (including dilatonic/electromagnetic polarizabilities) is essential for realistic compact objects and will enrich the phenomenology.
  \item \emph{Radiation and reaction.} Computing real emission (gravitational, electromagnetic, and dilatonic) and the accompanying radiation–reaction forces (in line with \cite{Caron-Huot:2023vxl}) will complete the bridge to waveform modeling for generic scattering and bound–state transitions, which is also work in progress \cite{arpan1}.
  \item \emph{Higher–curvature corrections.} Including $\alpha'$ \cite{Metsaev:1987zx}(higher–derivative) operators from the stringy effective action will test how UV-motivated terms imprint themselves on classical observables and whether symmetry protection persists at higher PM orders.
\end{itemize}

%\newpage
%\begin{figure}
%\centering
%\scalebox{0.3}{\includegraphics{Box_topo.png}}
%\caption{}\label{Box}
%\end{figure}
%\textbf{Penguin topologies:}
%\textbf{\begin{figure}
%\centering
%\scalebox{0.3}{\includegraphics{penguin.png}}
%\caption{}\label{penguin}
%\end{figure}}
%\newpage
